\documentclass[11pt]{article}
\usepackage[T1]{fontenc}
\usepackage[utf8]{inputenc}
\usepackage{lmodern}
\usepackage[margin=1in]{geometry}
\usepackage{amsmath,amssymb,amsthm,mathtools}
\usepackage{booktabs,tabularx,array}
\usepackage{microtype}
\usepackage{hyperref}
\hypersetup{colorlinks=true,linkcolor=blue,citecolor=blue,urlcolor=blue,
  pdftitle={Wilf's inequality for numerical semigroups with four generators},
  pdfauthor={Karthik Sethuraman}}
\newcolumntype{Y}{>{\raggedright\arraybackslash}X}
\newcommand{\Beta}{\mathrm{B}}
\numberwithin{equation}{section}
\newtheorem{theorem}{Theorem}[section]
\newtheorem{lemma}[theorem]{Lemma}
\newtheorem{corollary}[theorem]{Corollary}
\newtheorem{proposition}[theorem]{Proposition}
\theoremstyle{definition}
\newtheorem{finitelemma}[theorem]{Finite lemma}
\newtheorem{example}[theorem]{Example}
\theoremstyle{remark}
\newtheorem{remark}[theorem]{Remark}
\allowdisplaybreaks[2]
\setcounter{tocdepth}{2}
\title{Wilf's inequality for numerical semigroups\\with four generators}
\author{Karthik Sethuraman}
\date{September 17, 2026}

\begin{document}
\maketitle
\begin{abstract}
We present a computer-assisted argument for Wilf's inequality for numerical
semigroups with four minimal generators. Preferred factorizations of an
Ap\'ery set give a finite lower ideal in three exponent coordinates. We
seek a quantitative surplus over the weighted downset mean inequality,
large enough to absorb the exact correction in the Ap\'ery genus formula.
Residue collisions allow at most one full-support minimal exclusion.
We separate the resulting geometric cases, prove a central-box/three-horn
decomposition, and combine analytic moment inequalities with seven exact
finite assertions. The three-plane case is entirely analytic. Each retained
finite assertion is checked by two separate implementations; interval
certificates cover closed real parameter regions.
We also identify dimension-independent identities and a geometric obstruction
to directly extending the three-horn method beyond four generators.
\end{abstract}

\section{Introduction}\label{sec:introduction}

A numerical semigroup is a cofinite additive submonoid of
\(\mathbb N=\{0,1,2,\ldots\}\). Its embedding dimension \(e\) is the
cardinality of its minimal generating set, its multiplicity \(m\) is its
least positive element, and its conductor \(c\) is the least integer such
that \(c+\mathbb N\) is contained in the semigroup. Write
\(n=|S\cap[0,c)|\). Its genus is \(g=|\mathbb N\setminus S|\),
so \(n=c-g\). Wilf's question \cite{W} asks whether
\[
en\ge c
\]
for every numerical semigroup. Fr\"oberg, Gottlieb, and H\"aggkvist
\cite{FGH} established the inequality in embedding dimension at most
three. This paper addresses the next dimension.

\begin{theorem}\label{thm:main}
Every numerical semigroup with four minimal generators satisfies
\[
W_4(S):=4|S\cap[0,c)|-c\ge0.
\]
\end{theorem}

\subsection{Foundational work}

The proof builds on Zhai's \cite{Z} preferred factorizations of
Ap\'ery elements and weighted mean inequality for finite lower ideals.
Hellus, Rechenauer, and Waldi \cite{HRW} developed the monomial
initial-ideal and residue-lattice interpretation, including restrictions
on the supports of minimal exclusions. The classical Ap\'ery genus
formula supplies the exact link from
weighted moments to Wilf's inequality; see also \cite{Book}.

Two published computer-assisted results dispose of multiplicities through
19: Bruns, Garc\'ia-S\'anchez, O'Neill, and Wilburne
\cite[Theorem~4.3]{B}, and Kliem and Stump
\cite[Proposition~6.9, published version]{KS}.
The clique-tree ingredient in our geometric decomposition is classical
\cite{G}; its needed finite form is proved here. Connections with
three-dimensional L-shapes \cite{L,C} and other approaches to Wilf's
problem are discussed in Section~\ref{sec:related}.

\subsection{General approach}

Write \(S=\langle m,a_1,a_2,a_3\rangle\), minimally. Choose one
lexicographically preferred factorization of each Ap\'ery element by
\(a_1,a_2,a_3\). Their exponent vectors form a lower ideal
\(T\subseteq\mathbb N^3\) of cardinality \(m\), with labels distinct
modulo \(m\). Set
\[
A=\min_i a_i,\qquad w_i=a_i/A,\qquad
M=c+m-1,\qquad H=M/A,\qquad s=\sum_{x\in T}x.
\]
The genus formula gives
\[
mW_4(S)=A(3mH-4w\cdot s)-m(m-1).
\]
Thus the weighted mean baseline \(3mH-4w\cdot s\ge0\) does not suffice:
the proof needs a positive surplus that pays for the correction term.

We pursue this surplus through geometry and residue collisions. A minimal
excluded vector and its representative have disjoint supports, so there
is at most one minimal exclusion with all coordinates positive. An ideal
without such an exclusion is a central anchored box with at most three
monotone rectangular horns. An ideal with one is obtained from such a
completion by deleting a single upper orthant. This decomposition makes
additive extremal bounds accessible to finite, exact dynamic programming.

For small multiplicity, an analytic reduction places every hypothetical
counterexample inside an explicit bounded generator box. For large
multiplicity, the degree of the possible positive corner and the degree
of \(T\) give an exhaustive partition. The corner \((1,1,1)\) is handled
by a counting and projection argument. Other cases use local centroid
gains, axis witnesses, or interval domination of all real weights in a
compact failure region. A finite degree strip also yields a continuous
moment gap by averaging translated lattice meshes.

The weakest retained conclusion is
\(3mH-4w\cdot s\ge m-29/10\) for the branches not already settled
directly. For \(m\ge30\), the exact moment identity and \(A\ge m+1\)
then give \(W_4>-1\); integrality gives \(W_4\ge0\).

\subsection{Organization}

Section~\ref{sec:proof} gives the proof and its seven finite assertions.
Section~\ref{sec:analytic} supplies three longer analytic arguments integral
to that proof. Section~\ref{sec:verification} specifies the computational
artifact and its checks. The final two mathematical sections discuss
extensions beyond four generators and related work. We conclude with an
account of how the proof was constructed using AI and exact computation.

\tableofcontents
\section{Proof of the four-generator inequality}\label{sec:proof}

All vector comparisons are coordinatewise. A lower ideal contains every
nonnegative integer vector below each of its points. A positive corner
means a minimal excluded point with all three coordinates positive.

\subsection{Apéry representatives, moments, and collisions}\label{sec:p1}

The Apéry set
\[
\operatorname{Ap}(S,m)=\{w\in S:w-m\notin S\}
\]
consists of the least semigroup element in each residue modulo \(m\).
It has \(m\) elements, every factorization of which avoids \(m\).
Minimality of the generators gives \(0,a_1,a_2,a_3\in\operatorname{Ap}(S,m)\),
with the generator factorizations unique. Its largest element is
\[
M=c+m-1:
\]
every \(w-m\) is at most the largest gap \(c-1\), while \(c-1+m\) is Apéry.

For each Apéry element choose the lexicographically least exponent vector
\(x\in\mathbb N^3\) with label \(\lambda(x)=a\cdot x=w\).
Let \(T\) be the chosen vectors. This is the preferred-factorization
construction of Zhai, with the initial-ideal interpretation of
Hellus--Rechenauer--Waldi \cite{Z,HRW}.

To prove lower closure, take \(y\le x\in T\) and \(v=x-y\).
If \(\lambda(y)-m\in S\), then \(\lambda(x)-m\in S\), a contradiction.
Thus \(\lambda(y)\) is Apéry. A lexicographically earlier factorization
\(y'\) would give \(y'+v<_{\rm lex}x\) of the same label. Consequently
\(y\in T\). In particular,
\[
|T|=m,\quad 0,e_1,e_2,e_3\in T,\quad
\lambda:T\longrightarrow\operatorname{Ap}(S,m)
\ \text{is bijective},
\]
and the labels on \(T\) are distinct modulo \(m\).

Let \(s=\sum_Tx\), \(w=a/A\), \(H=M/A=\max_Tw\cdot x\), and
\(R=\max_T|x|_1\). Every \(w_i\ge1\), so \(R\le H\).
Writing the least representative of residue \(r\) as \(r+k_rm\)
and counting its \(k_r\) gaps proves the classical genus formula
(also \cite[Lemma~1]{Z})
\[
g=\frac1m\sum_{\operatorname{Ap}}w-\frac{m-1}{2}.
\]
Since \(W_4=3c-4g\), substitution gives the exact identity
\begin{equation}
\boxed{mW_4=A D_0-m(m-1),\qquad D_0=3mH-4w\cdot s.}
\label{eq:p:1}
\end{equation}

A minimal excluded point, or \emph{corner}, is \(p\notin T\) with
\(p-e_i\in T\) whenever \(p_i>0\). Let \(\rho(p)\in T\) have its residue.
If both \(p_i,\rho(p)_i>0\), subtracting \(e_i\) gives distinct included
points of equal residue, impossible. Likewise two distinct corners
sharing a positive coordinate cannot have the same residue.
Hence a positive corner has representative zero, and there is at most one
positive corner. All other corners have a zero coordinate.
These are the support/collision consequences of \cite[Proposition~2.6]{HRW}.

\subsection{Coordinate lines and two centroid bounds}\label{sec:p2}

For each of the three coordinate directions partition \(T\) into its
nonempty coordinate lines. Use an \emph{indexed} line family: intersecting
geometric lines in different directions remain distinct.
A line of length \(L_\ell\) starts at coordinate zero and has top \(t_\ell\).
Arithmetic progression sums give
\begin{equation}
\sum_\ell L_\ell=3m,\qquad
\sum_\ell L_\ell t_\ell=4s.
\label{eq:p:2}
\end{equation}
In a fixed direction its varying-coordinate contribution is twice that
coordinate's moment; the other coordinates contribute their moments once.
Thus, for any positive weights and any allowance \(Q\ge\max_Tw\cdot x\),
\begin{equation}
D_Q:=3mQ-4w\cdot s
=\sum_\ell L_\ell(Q-w\cdot t_\ell)\ge0.
\label{eq:p:3}
\end{equation}
This is the dimension-three weighted downset baseline of \cite[Lemma~3]{Z}.

Define
\[
\rho_2(t)=\max\{|v|_1:v\in\mathbb N^3,\ |v|_1\le2,\ t+v\in T\},
\qquad U_2=\sum_\ell L_\ell\rho_2(t_\ell).
\]
Choose a maximizing destination for each line. The total destination mass
is \(3m\), and its moment is \(4s+\beta\), with \(\beta\ge0\)
coordinatewise and \(|\beta|_1=U_2\). Each destination has weight at most
\(Q\). For weights at least one this proves
\begin{equation}
D_Q\ge U_2.
\label{eq:p:4}
\end{equation}
Only ten offsets are possible: zero, three units, three doubled units,
and three sums of distinct units.

More generally, any nonnegative point masses \(\alpha_x\) on \(T\)
with total \(3m\) and residual
\(\beta=\sum_x\alpha_xx-4s\ge0\) give \(D_Q\ge|\beta|_1\).
Indeed
\(D_Q=\sum_x\alpha_x(Q-w\cdot x)+w\cdot\beta\).
This identity also proves the explicit exceptional witnesses below.

If \(q_i=\max_Tx_i>0\), \(q_*=\max_iq_i\), put
\begin{equation}
G=q_*\left(3m-4\sum_i s_i/q_i\right).
\label{eq:p:5}
\end{equation}
When \(C=3m-4\sum_i s_i/q_i\ge0\), assign mass \(4s_i/q_i\)
to \(q_ie_i\), and mass \(C\) to a longest-axis endpoint.
The total is \(3m\), with residual \(Cq_*e_j\ge0\), proving \(D_Q\ge G\).
When \(C<0\), \eqref{eq:p:3} gives \(D_Q\ge0>G\).
In all cases \(D_Q\ge\max(U_2,G)\).

\subsection{The final-window projection inequality}\label{sec:p3}

Let \(\pi_j\) delete coordinate \(j\), and put
\[
P(T)=\sum_j|\pi_jT|,\quad
Z=\{x\in T:M-\lambda(x)<m\},\quad K=\operatorname{Max}(T).
\]
The set \(Z\) is nonempty. Its points are maximal, since an included
successor would increase their labels by \(a_j>m\) beyond \(M\).
For a nonempty point set \(X\), define
\[
E(X)=\sum_X|x|_1-\sum_j\max_Xx_j,\qquad
\Phi(T,X)=P(T)-3|X|-E(X).
\]
Here \(E(X)\ge0\). Adding a point increases each coordinate sum at least
as much as its coordinate maximum, so \(E(X)\le E(Y)\) for \(X\subseteq Y\).
Consequently \(\Phi(T,Z)\ge\Phi(T,K)\).

We prove
\begin{equation}
\boxed{mW_4\ge m\Phi(T,Z)+E(Z).}
\label{eq:p:6}
\end{equation}
Use the unnormalized weights \(a\) in the line formula, and write
\(M-\lambda(t_\ell)=m d_\ell+\beta_\ell\), \(0\le\beta_\ell<m\).
For a direction-\(j\) line of length \(L\), define
\[
R_j(L)=\sum_{r=0}^{L-1}[ra_j]_m,\quad
C_\ell=\#\{r<L:\beta_\ell+[ra_j]_m\ge m\},\quad
\sigma_\ell=mL d_\ell+mC_\ell-R_j(L).
\]
Reading the line backward gives
\[
L(M-\lambda(t_\ell))
=\sum_{x\in\ell}[M-\lambda(x)]_m+\sigma_\ell.
\]
Each direction partitions all \(m\) residues. Summing and using \eqref{eq:p:1} gives
\(mW_4=\binom m2+\sum_\ell\sigma_\ell\).

A deep line, whose top is outside \(Z\), has \(d_\ell\ge1\).
Its positive axis-prefix residues are distinct nonzero residues on \(T\);
therefore \(R_j(L)\le m(L-1)-\binom L2\) and \(\sigma_\ell\ge m\).
There are \(P(T)-3|Z|\) deep lines, because every \(z\in Z\) tops exactly
one line in each direction.

On a final line \(d_\ell=0\), so \(\sigma_\ell\ge-R_j(L)\).
Across final lines each positive axis-prefix point appears once, with
exactly \(E(Z)\) repeated occurrences: its multiplicity in direction \(j\)
is \(\#\{z\in Z:z_j\ge r\}\).
Distinct axis-prefix points have distinct nonzero residues.
The sum \(\binom m2\) pays for one copy of each used residue;
each repeated copy costs at most \(m-1\).
Thus
\[
mW_4\ge m(P(T)-3|Z|)-(m-1)E(Z),
\]
which is \eqref{eq:p:6}.
In particular \(\Phi(T,K)\ge0\) suffices to prove Wilf.

\subsection{Phase and rectangular thickening}\label{sec:p4}

Normalize by attained height: \(u_i=w_i/H\), \(v=\min_i u_i\),
\(U=\max_i u_i\), \(\sigma=\sum_i u_i\),
\(\mu_i=\mathbb E_T(u_iX_i)\), and
\(\kappa=3-4\sum_i\mu_i=D_0/(mH)\).
Units give \(u_i\le1\). The mean slack of direction \(j\) line tops is
defined as follows. For \(x\in T\), let \(t_j(x)\) be the top of the
direction-\(j\) line through \(x\), and put
\[
\delta_j(x)=1-u\cdot t_j(x),\qquad
\bar\delta_j=\frac1m\sum_{x\in T}\delta_j(x).
\]
Thus \(\bar\delta_j\) is the length-weighted mean over direction-\(j\)
lines, and the line moment identities give
\[
\bar\delta_j=(1+\kappa)/4-\mu_j,\qquad
\sum_j\bar\delta_j=\kappa\ge0.
\]

Choose coordinate \(k\) of weight \(v\).
Every interval of length \(h\) has sawtooth integral
\[
\int_I[r]_q\,dr\ge\frac{qh^2}{2(h+q)}.
\]
To prove this, write \(h=nq+r_0\): \(n\) full periods contribute
\(nq^2/2\), the remaining arc at least \(r_0^2/2\);
Cauchy--Schwarz gives \((nq^2+r_0^2)/2\ge h^2/(2(n+1))\),
and \(n+1\le(h+q)/q\).

On a coordinate-\(k\) line of length \(L\), write \(\delta_j(t)\) for the
slack at its point with \(k\)-coordinate \(t\). For \(j\ne k\), these
slacks have residues \([\theta-tv]_{u_j}\).
Nonnegative slack and
\([\theta-tv+r]_{u_j}\le\delta_j(t)+r\), \(0\le r\le v\),
bound the integral over the concatenated intervals of total length \(Lv\).
Since \(Lv\le1+v\), division gives
\[
\operatorname{mean}_{\rm line}\delta_j
\ge\frac{u_jvL}{2(1+v+u_j)}-\frac v2.
\]
Length-weighted averaging uses \(v\sum L^2/m=2\mu_k+v\), hence
\[
2\bar\delta_j(1+v+u_j)\ge2u_j\mu_k-v(1+v).
\]
Sum over \(j\ne k\), substitute \(\mu_k=(1+\kappa)/4-\bar\delta_k\),
and use \(U\le\sigma-2v\).
The remaining slack coefficient is the middle weight minus \(1+v\),
nonpositive because that weight is at most one. Discarding it gives
\begin{equation}
\boxed{\sigma(1-3\kappa)\le4\kappa+5v-3\kappa v+4v^2.}
\label{eq:p:7}
\end{equation}
If \(H>3\) and \(D_0<m\), then \(0\le\kappa<v=1/H<1/3\).
The right side divided by \(1-3\kappa\) is strictly increasing in \(\kappa\),
with derivative numerator \(4+12v+12v^2>0\).
Writing \(B=\sum_iw_i\) proves the necessary failure cutoff
\begin{equation}
\boxed{B<9+28/(H-3).}
\label{eq:p:8}
\end{equation}

Thicken \(T\) into equal-volume half-open boxes
\[
K_T=\bigcup_{x\in T}\prod_i
\left[\frac{w_ix_i}{H+B},\frac{w_i(x_i+1)}{H+B}\right).
\]
It is a positive-volume downset in the unit simplex.
Its mean coordinate sum is
\((w\cdot s/m+B/2)/(H+B)\), giving
\begin{equation}
\boxed{\kappa_c(K_T):=3-4\mathbb E_{K_T}|Y|_1
=\frac{D_0/m+B}{H+B}.}
\label{eq:p:9}
\end{equation}
The complement is the union of upper orthants with vertices
\((w_ip_i/(H+B))_i\), for the discrete minimal exclusions \(p\).
Positive diagonal scaling preserves their supports and minimality.
The same calculation with an allowance \(Q\) in place of \(H\) replaces
\(D_0\) by \(D_Q\) in the identity.
A continuous lower bound \(\kappa_c\ge\gamma\) therefore gives
\begin{equation}
D_0/m\ge\gamma H-(1-\gamma)B,\qquad
D_0<m\ \Longrightarrow\ \gamma H<1+(1-\gamma)B.
\label{eq:p:10}
\end{equation}

\begin{lemma}[Height cutoff]\label{lemma:height-cutoff}
Suppose \(H>3\), \(D_0<m\), and
\[
\frac{D_0}{m}\ge\alpha H-\beta B,\qquad
\alpha>0,\quad 0\le\beta\le3.
\]
Then
\begin{equation}
\boxed{H<C:=3+\frac{1+9\beta}{\alpha}}.
\label{eq:height-cutoff}
\end{equation}
\end{lemma}
\begin{proof}
The assumed moment bound gives \(\alpha H<1+\beta B\).
Substitute \eqref{eq:p:8} and multiply by the positive quantity \(H-3\):
\[
\alpha H^2-(3\alpha+1+9\beta)H+3-\beta
=\alpha H(H-C)+3-\beta<0.
\]
If \(H\ge C\), the left side is nonnegative because \(\alpha>0\)
and \(3-\beta\ge0\), a contradiction.
\end{proof}

\begin{table}[htbp]
\centering
\begin{tabular}{@{}lccc@{}}
\toprule
Branch & \(\alpha\) & \(\beta\) & Bound \\
\midrule
No full-support corner & \(1/3\) & \(2/3\) & \(H<24\) \\
Short corner, high height & \(1/3\) & \(4/3\) & \(H<42\) \\
Residual high height & \(5/42\) & \(37/42\) & \(H<78\) \\
\bottomrule
\end{tabular}
\caption{Parameters for the common height-cutoff argument. The interval
certificates cover the corresponding closed supersets.}
\label{tab:height-cutoffs}
\end{table}

\subsection{Structural geometry and an exact interval method}\label{sec:p5}

Theorem~\ref{theorem:a:1.1}, proved in Section~\ref{sec:analytic-a},
proves that every finite no-full-corner lower ideal \(U\) is a disjoint
central anchored box \([0,z]\) and at most three outward horns.
Horn \(i\) consists of rectangles at levels \(t>z_i\), with transverse
caps at most \(z_j,z_k\), nonincreasing with \(t\).
The proof constructs a chordal compatibility graph, proves its needed
clique-tree property, and uses the median of three coordinate maximizers.

For a unique positive corner \(p\), deleting its generator from the
upper-ideal complement gives a finite no-full-corner \(U\) with
\begin{equation}
T=U\setminus(p+\mathbb N^3).
\label{eq:p:11}
\end{equation}
Pure-axis exclusions remain, so \(U\) is finite and has the same axes.
Clipping the center and horns preserves their disjointness.

We use the following exact upper-bound method in the three interval checks.
For these checks, permute coordinates and weights together so that
\(w=(1,b,d)\), with \(1\le b\le d\).
At scale \(q=4096\), integer endpoints \(L=(B_0,C_0,Q_0)\),
\(V=(B_1,C_1,Q_1)\) enclose a closed real parameter box for \((b,d,Q)\).
Put
\[
\ell=(q,B_0,C_0),\quad u=(q,B_1,C_1),\quad
N_i=\lfloor Q_1/\ell_i\rfloor,\quad
\psi(x)=4u\cdot x+q-3Q_0.
\]
Every relevant retained point satisfies \(\ell\cdot x\le Q_1\).
Every axis of \(U\) survives clipping, so \(U\subseteq\prod_i[0,N_i]\).
Pointwise monotonicity gives, for every real parameter in the box,
\begin{equation}
q\bigl(m-(3mQ-4w\cdot s)\bigr)\le\sum_T\psi(x).
\label{eq:p:12}
\end{equation}

For an inclusive box \([a,d]\), its count is \(\prod_i(d_i-a_i+1)\)
and twice its objective moment is that count times \(\sum_i u_i(a_i+d_i)\).
For a clipped box subtract the upper box \([\max(a,p),d]\).
The retained feasibility maximum is
\begin{equation}
\max_{\substack{i\\a_i\le\min(d_i,p_i-1)}}
\left(\ell_i\min(d_i,p_i-1)+\sum_{j\ne i}\ell_jd_j\right).
\label{eq:p:13}
\end{equation}
This maximizes over the parts with some coordinate below \(p\).
An independent formula partitions by the \emph{first} such coordinate,
summing disjoint boxes with earlier coordinates at least their \(p_j\).
Their maximum heights and summed moments give the same statistics.
Without clipping use the full box and its upper-corner height.

Let \(C_i(t,r,s)\) be the exact section score, and call it feasible
when its retained maximum height is at most \(Q_1\).
For a possibly empty nested horn use
Lemma~\ref{lemma:horn-recurrence}, terminating beyond its axis cap.
Enumerate every feasible retained central box and add
its score to \(\sum_iF_i(z_i+1;z_j,z_k)\).
The structural theorem then proves an upper bound \(V_p(L,V)\) on \eqref{eq:p:12}.
Independent horns may enlarge the class; they never decrease its upper bound.

Order tightening replaces endpoints by
\begin{equation}
L^\sharp=(B_0,\max(B_0,C_0),\max(B_0,C_0,Q_0)),\quad
V^\sharp=(\min(B_1,C_1,Q_1),\min(C_1,Q_1),Q_1).
\label{eq:p:15}
\end{equation}
Its ordered intersection is unchanged; inversion proves emptiness.
Splitting a tightened box at an interior integer wall yields two closed
children sharing that wall and covering the parent.
A finite tree reaching every node exactly once and ending only at checked
empty, analytic, or exact-DP leaves covers all real root parameters.


\subsection{Small multiplicity}\label{sec:p6}

The two published computer-assisted results \cite[Theorem~4.3]{B} and
\cite[Proposition~6.9]{KS} settle all numerical semigroups of multiplicity at
most 18 and exactly 19, respectively, without an embedding-dimension
restriction.

For \(20\le m\le29\) we need an analytic counterexample reduction.
Theorem~\ref{theorem:b:1.1}, proved in Section~\ref{sec:analytic-b},
proves Wilf whenever \(T=\{x:a\cdot x\le M\}\), including an analytic
residue obstruction to its formal six-column case.

If the ideal is not full weighted, choose an omitted point of weight
at most \(M\) of least degree. It is a minimal excluded point \(p\)
with \(a\cdot p\le M\). Its direction-\(i\) predecessor line, for \(p_i>0\),
has top \(p-e_i\) and length exactly \(p_i\).
Their contribution to the unnormalized line slack is
\[
\sum_{i:p_i>0}p_i(M-a\cdot p+a_i)
=M+(|p|_1-1)(M-a\cdot p)\ge M.
\]
All other slacks are nonnegative. Therefore \(W_4<0\), which implies
the integer bound \(W_4\le-1\), forces
\begin{equation}
\boxed{M\le3mM-4a\cdot s\le m(m-2).}
\label{eq:p:16}
\end{equation}
Every generator \(a_i\) is Apéry and at most \(M\).
After sorting, any counterexample thus belongs to the \emph{inclusive} box
\[
m<a_1<a_2<a_3\le B_m:=m(m-2).
\]

\begin{finitelemma}[Generator-box]\label{finite:generator-box}
For each \(20\le m\le29\), every increasing
triple in this box with \(\gcd(m,a_1,a_2,a_3)=1\) and all four generators
minimal has either largest Apéry element \(M>B_m\), or \(W_4\ge0\).

\end{finitelemma}

Every raw triple is classified by successive exact tests:
\(a_1\in\langle m\rangle\), \(a_2\in\langle m,a_1\rangle\),
\(a_3\in\langle m,a_1,a_2\rangle\), then a nontrivial gcd, or validity.
Larger generators cannot make an earlier one redundant.
One path computes exact residue shortest distances \(d_r\), adjoining
a generator by two forward passes around each residue cycle.
A shortest extension uses fewer than the cycle length many new copies;
two passes propagate every possible starting residue around the cycle.
The second path uses increasing ordinary membership:
\(I_j(n)=I_{j-1}(n)\vee(n\ge a_j\ \&\ I_j(n-a_j))\).
It tests the complete window \([B_m-m+1,B_m]\), which is full exactly
when \(M\le B_m\). A full window and closure under \(+m\) prove every
larger integer is included, so gaps and conductor are fully determined.

The generator-box finite lemma contradicts \eqref{eq:p:16} for any negative example.
Thus all multiplicities through 29 are settled.

\subsection{Large multiplicity: exhaustive routing and the three-plane case}\label{sec:p7}

Now \(m\ge30\). If there is no positive corner use Section~\ref{sec:p8}.
Otherwise let \(p\) be the unique positive corner.
If \(|p|_1=3\), then \(p=(1,1,1)\) and the following analytic argument applies.
Degree four gives a permutation of \((2,1,1)\), handled in Section~\ref{sec:p9}.
For \(|p|_1\ge5\), degree \(R\le6\) uses Section~\ref{sec:p10}; otherwise \(R\ge7\)
uses Section~\ref{sec:p11}. These are disjoint exhaustive alternatives.
Minimality gives \(|p|_1-1\le R\); consequently Section~\ref{sec:p10} has
\(5\le|p|_1\le7\), while Section~\ref{sec:p11} has \(H\ge R\ge7\).
All coordinate permutations are applied simultaneously to weights.

Suppose \(p=(1,1,1)\). Every point of \(T\) lies in a coordinate plane.
A mixed corner \(q\) in the \(ij\)-plane has representative \(\gamma e_k\),
\(1\le\gamma\le h_k=\max_Tx_k\). The included point
\(q-e_i-e_j\) has residue \((\gamma+1)a_k\), since \(a\cdot p\equiv0\).
If \(\gamma<h_k\) it collides with \((\gamma+1)e_k\).
Thus all mixed corners in that plane have representative \(h_ke_k\);
the collision rule permits at most one.

A geometric count proves \(\Phi(T,K)\ge0\) for \(m\ge30\).
Write \(h=h_1+h_2+h_3\ge3\).
Each plane is its full rectangle or that rectangle with one upper-right
quadrant deleted; it contains \((1,1)\).
Its frontier has one of the following forms, allowing transpose:

\begin{center}
\small
\begin{tabular}{@{}lllll@{}}
\toprule
Type & Mixed maxima & Mixed point count \(c_{ij}\) & Correction \(\beta_{ij}\) & Axis charge \(\nu_{ij}\) \\
\midrule
Rectangle & \((h_i,h_j)\) & \(h_ih_j\) & 0 & 0 \\
Split, \(1\le b<h_i,\ 1\le d<h_j\) & \((b,h_j),(h_i,d)\) & \(bh_j+dh_i-bd\) & \(b+d\) & 0 \\
Side cap at \(i\), \(1\le b<h_i\) & \((b,h_j)\) & \(bh_j\) & \(b\) & \(h_i\) \\
\bottomrule
\end{tabular}
\end{center}

Let \(s_f\) count split planes, \(I\) the globally maximal axis directions,
\(J=\sum_{i\in I}h_i\), \(\Beta=\sum\beta_{ij}\),
\(C=\sum(c_{ij}-\beta_{ij})\), and \(N=\sum\nu_{ij}\).
Every split plane requires both incident axis caps to be at least two;
checking \(s_f=0,1,2,3\) gives \(h\ge s_f+3\).
Every mixed planar maximum is globally maximal. An axis top is globally
maximal exactly when both incident planes side-cap that axis.
Each side-cap plane can cap only one axis, while a split plane caps none.
Consequently \(2|I|\le3-s_f\), so \(s_f+|I|\le3\) and
\(k=|K|=3+s_f+|I|\le6\).
For a split,
\(c_{ij}-\beta_{ij}
=bd+b(h_j-d-1)+d(h_i-b-1)\ge b+d-1\).
For a side cap \(c_{ij}-\beta_{ij}\ge0\) and \(\beta_{ij}\le\nu_{ij}\).
Hence
\[
\Beta\le C+s_f+N,\qquad N\ge2J.
\]
Direct origin/axis/plane counting yields
\[
m=1+h+\Beta+C,\quad P(T)=m+h+2,\quad E(K)=h+\Beta-N+J.
\]
Set \(Q_*=h+C+N-J\). Then \(\Phi(T,K)=Q_*+3-3k\).
If this were negative, integrality would give \(Q_*\le3k-4\le14\), and
\[
m=1+Q_*+\Beta-N+J
\le1+2Q_*-h+s_f-N+2J
\le2Q_*-2\le26.
\]
This contradicts \(m\ge30\). Therefore \eqref{eq:p:6} proves \(W_4\ge0\).

\subsection{No full-support corner}\label{sec:p8}

Theorem~\ref{theorem:c:5.2}, proved in Section~\ref{sec:analytic-c},
proves \(\kappa_c\ge1/3\) for every positive-volume finite-step
central-box/horn set in the unit simplex.
That section proves that degree at most four in this discrete class
implies cardinality at most 29, by eleven explicit central-cap bounds.

If \(D_0<m-1\) at \(m\ge30\), then \(H\ge5\), since \(R\le H<5\)
would force degree at most four. Sort weights to \((1,b,d)\), so
\(1\le b\le d\le H\). The continuous bound and \eqref{eq:p:10}
give the first row of Table~\ref{tab:height-cutoffs}, so
Lemma~\ref{lemma:height-cutoff} gives \(H<24\).
For any allowance \(Q\), the continuous theorem also gives
\(D_Q/m\ge(Q-2B)/3\).

\begin{finitelemma}[No-corner interval]\label{finite:no-corner-interval}
For all real
\(1\le b\le d\le Q\), \(5\le Q\le24\), and every nonempty finite
no-full-corner ideal of attained height at most \(Q\),
\[
m-D_Q\le1.
\]

\end{finitelemma}

This is stronger than the genuine branch domain: no units, residue
labels, attained allowance, or cardinality restriction is required.
Use the method in Section~\ref{sec:p5}, with no clipping.
A DP leaf requires its exact bound at most \(q\).
An analytic leaf uses
\(Q_0\ge2(q+B_1+C_1)+3q\), which implies \(D_Q\ge m\).
The interval certificate covers the closed root from height 5 to 24.
The finite lemma contradicts the hypothetical failure.
Thus \(D_0\ge m-1\) throughout this branch.

\subsection{The short corner}\label{sec:p9}

Orient \(p=(2,1,1)\), without prescribing the least weight's direction.
For any degree allowance \(r\ge3\),
\begin{equation}
|T|\le2r^2-r+3.
\label{eq:p:17}
\end{equation}
Indeed the degree-\(r\) simplex minus the \(p\)-orthant has \(2r^2+2\)
points. For each \(1\le j\le r-1\), the excluded point \((1,j,r-j)\)
cannot dominate the unique full-support corner.
It must dominate a planar corner; at least one of
\((0,j,r-j),(1,0,r-j),(1,j,0)\) is therefore omitted.
These \(r-1\) triples are pairwise disjoint, forcing \eqref{eq:p:17}.
In particular \(H<6\) gives \(R\le5\) and \(30\le m\le48\).
The endpoint \(H=6\) belongs to the high-height case.

\subsubsection{High height}

For a translated planar lower ideal of offset \(h\) in height \(H\),
the dimension-two line identity gives mean weight at most \((2H+h)/3\).
Partition \(T\) into \(y=0\); \(y\ge1,z=0\);
\(x=0,y,z\ge1\); and \(x=1,y,z\ge1\).
These are translated planar lower ideals, with offsets respectively
\(0,w_2,w_2+w_3,B\). Thus \(D_0/m\ge(H-4B)/3\).
If \(H\ge6\) and \(D_0<m-1\), Lemma~\ref{lemma:height-cutoff}
with the second row of Table~\ref{tab:height-cutoffs} gives \(H<42\).
We retain the closed height-42 superset.

\begin{finitelemma}[Short-corner high]\label{finite:short-corner-high}
For all
\(1\le b\le d\le Q\), \(6\le Q\le42\), each coordinate permutation of
\(p=(2,1,1)\), and every nonempty finite no-full-corner \(U\),
the clipped \(T=U\setminus(p+\mathbb N^3)\) of height at most \(Q\)
satisfies \(m-D_Q\le1\).

\end{finitelemma}

Apply Section~\ref{sec:p5}'s retained-height feasibility and clipped scores to all
three corners. A DP leaf requires every corner bound at most \(q\).
The planar rule \(Q_0\ge4(q+B_1+C_1)+3q\) is also sound, including equality.
The complete closed tree establishes the finite lemma, contradicting failure.

\subsubsection{Low height: exact profiles and two inline exceptions}

A mixed \(xy\) corner has representative \(\gamma e_3\ne0\).
If its \(x\)-coordinate is at least two, subtracting \(2e_1+e_2\)
gives an included point with residue \((\gamma+1)a_3\).
It forces \(\gamma=h_3\); at most one such corner exists.
The antichain permits at most one additional corner at \(x=1\).
The \(xz\) plane is identical. A \(yz\) corner yields residue
\((\gamma+2)a_1\) on subtracting \(e_2+e_3\), forcing
\(\gamma\in\{h_1-1,h_1\}\).
Hence each plane has at most two mixed corners, with the refined incident
plane restriction, and contains respectively \((2,1),(2,1),(1,1)\).

A degree-five plane profile is \(f=(f_0,\ldots,f_5)\),
\(0\le f_i\le6-i\), nonincreasing, encoding \(\{(i,j):j<f_i\}\).
Let \(\ell(f)\) count positive columns and
\(\Delta(f)=\{i:1\le i<\ell(f),\ f_i<f_{i-1}\}\).
These, and only these, are mixed corners; a drop to zero is a pure-axis corner.

For \(xy,xz\), require \(f_2\ge2\), \(|\Delta(f)|\le2\), and at most
one drop at index at least two. For \(yz\), require \(h_1\ge2\)
and \(|\Delta(h)|\le2\).
Use all compatible ordered triples \(f,g,h\), with
\[
\ell(f)=\ell(g),\quad f_0=\ell(h),\quad g_0=h_0,
\]
and reconstruct
\begin{equation}
T=\{(x,y,z):y<f_x,\ z<g_x,\ z<h_y,\ (x,y,z)\not\ge(2,1,1)\}.
\label{eq:p:18}
\end{equation}
Keep exactly those of degree at most five and \(30\le|T|\le48\).
Matching axes recover each plane section exactly.
The required plane entries include all three predecessors of \(p\).
Any extra reconstructed point in a genuine ideal would dominate a planar
minimal exclusion, violating its corresponding pair test.
Thus every genuine low-height ideal is recovered uniquely.
Profiles alone do not guarantee total degree: the literal degree check is essential.

\begin{finitelemma}[Short-corner low]\label{finite:short-corner-low}
Every such triple satisfies
\(U_2\ge m-1\), except possibly the two keys
\[
\begin{aligned}
E_0&=((4,3,3,2,0,0),(6,3,3,2,0,0),(6,4,2,2,0,0)),\\
E_1&=((6,3,3,2,0,0),(4,3,3,2,0,0),(4,4,2,2,1,1)).
\end{aligned}
\]

\end{finitelemma}

For \(E_0\), direct reconstruction gives \(m=30\), \(s=(27,25,34)\).

The included points \((0,0,5),(2,2,0),(3,1,0)\) with masses \(38,48,4\)
have total mass 90 and residual
\((108,100,190)-4s=(0,0,54)\).
Section~\ref{sec:p2} therefore gives \(D_Q\ge54\ge29\) for all admissible real weights.
The second ideal and witness are obtained by interchanging \(y,z\),
including transposition of the \(yz\) profile.

Together with \eqref{eq:p:4}, the finite lemma proves \(D_0\ge m-1\) in low height.
This completes the short-corner branch.

\subsection{Residual degree at most six}\label{sec:p10}

Let \(5\le P=|p|_1\le7\), \(R\le6\).
In an \(ij\)-plane, distinct mixed corners have distinct representatives
\(\gamma e_k\ne0\). At most \(p_i-1\) have coordinate \(q_i<p_i\),
and at most \(p_j-1\) have \(q_j<p_j\), by antichain distinct coordinates.
For the remaining corners, subtract \(p_ie_i+p_je_j\).
The included result has residue \((\gamma+p_k)a_k\);
collision with the \(k\)-axis forces \(\gamma>h_k-p_k\).
At most \(p_k\) distinct representatives remain. Each plane thus has at most
\begin{equation}
(p_i-1)+(p_j-1)+p_k=P-2
\label{eq:p:19}
\end{equation}
mixed corners.

Define \(F_i=\{x\in T:x+e_i\notin T\}\), \(n_k=1+\max_Tx_k\).
For every residue-bijective lower ideal,
\begin{equation}
|F_i\cap F_j|\le2n_k.
\label{eq:p:20}
\end{equation}
To prove it, let \(C_{ij}\) contain excluded points with both positive
\(i,j\) coordinates and both respective predecessors included.
In each nonempty fixed-\(k\) slice, a planar ideal with \(h\) maxima has
\(h-1\) mixed corners, by its constant positive-height blocks.
All \(n_k\) slices are nonempty, hence
\(|C_{ij}|=|F_i\cap F_j|-n_k\).
Subtracting an included predecessor forces the representative of
each \(q\in C_{ij}\) onto the \(k\)-axis, and distinct \(q\) have distinct
representatives by comparing their \(i\)-predecessors.
This injects \(C_{ij}\) into the \(n_k\) axis points, proving \eqref{eq:p:20}.
Slice corners need not be globally minimal; only these two predecessors are used.

Sort coordinates of \(p\), simultaneously permuting weights.
The nine corners are
\[
(1,1,3),(1,2,2),(1,1,4),(1,2,3),(2,2,2),
(1,1,5),(1,2,4),(1,3,3),(2,2,3).
\]
Use all nonincreasing seven-entry profiles \(0\le f_i\le7-i\).
In an \(ij\)-plane require \(f_{p_i}\ge p_j+1\) and at most \(P-2\)
positive height drops. Impose the three shared-axis compatibilities
as in \eqref{eq:p:18}, and reconstruct with its own \(p\)-orthant deleted.
Keep precisely degree at most six, cardinality at least 30,
and the three bounds \eqref{eq:p:20}. No maxima, realizability, modular, or erosion filter is used.
The same plane-recovery and excluded-corner argument as for \eqref{eq:p:18}
proves complete, unique coverage of genuine oriented ideals.

\begin{finitelemma}[Degree-six]\label{finite:degree-six}
On this complete geometric family,
\begin{equation}
\max(U_2,G)\ge m.
\label{eq:p:21}
\end{equation}

\end{finitelemma}

The local integer predicate and the denominator-cleared axis predicate
are exact. With \(d=q_1q_2q_3>0\), the latter is
\[
q_*\left(3md-4\sum_i s_i\prod_{j\ne i}q_j\right)\ge md.
\]
Section~\ref{sec:p2} and \eqref{eq:p:21} give \(D_0\ge m\) for this branch.

\subsection{Higher degree: a finite strip and a real interval bound}\label{sec:p11}

Assume \(R\ge7\), \(|p|_1\ge5\). We prove \(D_0\ge m-29/10\).
The first finite premise supplies a continuous inequality; that inequality
restricts a hypothetical failure to a compact real parameter region.
The second finite premise excludes the entire region.

\subsubsection{The three-allowance strip}

\begin{finitelemma}[Strip]\label{finite:strip}
For \(r\in\{18,19,20\}\), every finite lower
ideal \(T'\subseteq\mathbb N^3\) of degree at most \(r\), with at most
one positive minimal exclusion, satisfies
\begin{equation}
3r|T'|-4\sum_{x\in T'}|x|_1\ge4|T'|.
\label{eq:p:22}
\end{equation}
Empty ideals and missing coordinate units are allowed.

\end{finitelemma}

The empty ideal satisfies the inequality directly. For nonempty ideals,
the following finite reduction is exhaustive. If there is a positive minimal
corner \(p'\), its predecessors imply \(|p'|_1\le r+1\).
Remove that exclusion generator to obtain a finite no-corner completion:
pure-axis exclusions remain, and deleting a positive orthant leaves every
axis point unchanged. Thus its axis caps are at most \(r\).
If there is no positive corner, take \(p'=(1,1,r-1)\); its orthant is
inactive on the degree-\(r\) simplex.
By symmetry enumerate exactly
\[
1\le p'_1\le p'_2\le p'_3,\qquad |p'|_1\le r+1.
\]
No minimality filter is imposed on this larger family.

Apply the decomposition of Section~\ref{sec:analytic-a}, clipped by \(p'\), to the score
\(f_r(x)=4|x|_1+4-3r\). All center and transverse caps range from zero
to \(r\). A section or center is feasible when its \textbf{retained}
maximum degree is at most \(r\), not when its removed upper corner is.
Apply Lemma~\ref{lemma:horn-recurrence} with the section sum of this
score, retained-degree feasibility, and termination beyond the degree
allowance.
The maximum of center score plus its three horn bounds is an upper bound
for \(\sum_{T'}f_r=4|T'|-D_r(T')\). Every configuration must have this
bound at most zero. This is a search over corner/allowance configurations
and exact dynamic-programming states.
Singletons, zero caps, and inactive corners are included.

\subsubsection{Transfer to the continuous gap}

Let \(K\) be a positive-volume downset in the unit simplex whose complement
is a finite union of upper orthants with at most one positive minimal vertex.
Put \(h=1/21\), and for \(z\in[0,h)^3\) define
\[
T_z=\{y\in\mathbb N^3:z+hy\in K\},\quad
r_z=\left\lfloor(1-|z|_1)/h\right\rfloor,\quad
\epsilon_z=1-|z|_1-hr_z.
\]
For \(z\ne0\), \(0<|z|_1<3h\), so
\(r_z\in\{18,19,20\}\) and \(0\le\epsilon_z<h\).
The zero translation is a null set and is omitted throughout the integral;
no allowance-21 assertion is needed.
Each excluded vertex \(g\) induces the discrete vertex
\[
g'_i=\max(0,\lceil(g_i-z_i)/h\rceil).
\]
A zero coordinate stays zero. Consequently at most one induced minimal
vertex is positive, even after redundant vertices are removed.
For almost every translation boundary conventions agree; the exceptional
translations lie in finitely many coordinate hyperplanes modulo \(h\).
The simplex condition gives \(|y|_1\le r_z\) for every \(y\in T_z\).
The hypotheses of \eqref{eq:p:22} therefore hold for \(T_z\), including when
units are absent. Expansion gives
\[
\sum_{y\in T_z}(3-4|z+hy|_1)
 =hD_{r_z}(T_z)+(3\epsilon_z-|z|_1)|T_z|
 \ge(4h-|z|_1)|T_z|.
\]
Integrating the partition \(x=z+hy\), whose Jacobian is one, gives
\[
3-4\mathbb E_K|X|_1
\ge4h-\sum_i\mathbb E_K(X_i\bmod h).
\]
A coordinate fiber of a downset is an interval \([0,L]\) up to endpoints.
Write \(L=qh+t\), \(0\le t<h\). Since
\[
\int_0^L(x\bmod h)\,dx=(qh^2+t^2)/2\le hL/2,
\]
Fubini implies each remainder average is at most \(h/2\). Hence
\begin{equation}
\boxed{\kappa_c(K)\ge5/42.}
\label{eq:p:24}
\end{equation}
The averaged remainder estimate is essential; a pointwise estimate alone
does not give this constant. The rectangular thickening from Section~\ref{sec:p4}
preserves excluded supports and meets every hypothesis above.

\subsubsection{Necessary region for a failure}

Sort the normalized weights to \(w=(1,b,d)\), and write \(B=1+b+d\).
The predecessors of the positive corner contain the units, so
\(1\le b\le d\le H\), and \(H\ge R\ge7\).
If \(D_0<m-29/10<m\), Section~\ref{sec:p4} and \eqref{eq:p:24} imply
\[
B<9+\frac{28}{H-3},\qquad
\frac{D_0}{m}\ge\frac{5H-37B}{42},\qquad 5H<42+37B.
\]
Lemma~\ref{lemma:height-cutoff}, with the final row of
Table~\ref{tab:height-cutoffs}, gives \(H<78\).
Also \(p-e_1\in T\) gives \(w\cdot p\le H+1\).
It suffices to cover the closed necessary failure region
\begin{equation}
\mathcal F=\{(b,d,Q):
1\le b\le d\le Q,\ 7\le Q\le78,\
1+b+d\le9+28/(Q-3),\
5Q\le42+37(1+b+d)\}.
\label{eq:p:25}
\end{equation}

\begin{finitelemma}[High-height]\label{finite:high-height}
For every \((b,d,Q)\in\mathcal F\),
every positive integer \(p\) with \(|p|_1\ge5\) and
\((1,b,d)\cdot p\le Q+1\), and every nonempty finite no-corner lower ideal \(U\),
put \(T'=U\setminus(p+\mathbb N^3)\). If
\(\max_{T'}(1,b,d)\cdot x\le Q\), then
\begin{equation}
|T'|-\left(3Q|T'|-4(1,b,d)\cdot\sum_{T'}x\right)\le29/10.
\label{eq:p:26}
\end{equation}

\end{finitelemma}

No residue, cardinality, minimality, unit, or surface restriction is used
in this enlarged geometric family.
Use the following interval tightening. At scale \(q=4096\) write endpoints
\(L=(B_0,C_0,H_0)\), \(U=(B_1,C_1,H_1)\) for \((b,d,Q)\).
First intersect with order:
\[
L\leftarrow(B_0,\max(B_0,C_0),\max(B_0,C_0,H_0)),\quad
U\leftarrow(\min(B_1,C_1,H_1),\min(C_1,H_1),H_1).
\]
If not inverted, compute, in this order,
\begin{equation}
\begin{aligned}
S_{\max}&=9q+\left\lceil28q^2/(H_0-3q)\right\rceil,\\
B_1&\leftarrow\min(B_1,\lceil(S_{\max}-q)/2\rceil),\\
C_1&\leftarrow\min(C_1,S_{\max}-q-B_0),\\
S_{\mathrm{upper}}&=\min(q+B_1+C_1,S_{\max}),\\
H_1&\leftarrow\min(H_1,\lceil(42q+37S_{\mathrm{upper}})/5\rceil).
\end{aligned}
\label{eq:p:27}
\end{equation}
Perform at most twenty iterations, stopping on inversion or unchanged
endpoints. Every upper rounding is outward. The inequalities in \eqref{eq:p:25},
the lower bound \(qQ\ge H_0\), and \(b\le d\) prove that each step retains
every relevant real point. Convergence is not required for soundness.

For a nonempty tightened box use Section~\ref{sec:p5} with
\(\ell=(q,B_0,C_0)\), \(u=(q,B_1,C_1)\), and
\(\psi(x)=4u\cdot x+q-3H_0\). Enumerate the \textbf{full Cartesian} corner list
\[
p_i\ge1,\quad5\le|p|_1\le\lfloor H_1/q\rfloor+1,\quad
\ell\cdot p\le H_1+q.
\]
This contains every relevant corner, since \(\ell_i\ge q\).
An empty list is accepted only as vacuous.
For every listed corner the retained-center/three-horn bound \(V_p\)
must satisfy \(10V_p\le29q\).
The closed root is \(((q,q,7q),(78q,78q,78q))\).
Each child is derived from its tightened parent; each split is strictly
interior, and its two children share the split wall. Checked empty leaves
and checked dynamic-programming leaves exhaust all nodes.
Induction proves \eqref{eq:p:26} throughout the closed real parameter region.
At \(Q=H\), it contradicts the proposed failure and proves this branch.

\subsection{Assembly of the proof}\label{sec:p12}

\begin{proof}[Proof of Theorem~\ref{thm:main}]

The branches are exhaustive. Multiplicity at most 29 is settled by Section~\ref{sec:p6}.
For \(m\ge30\), Section~\ref{sec:p7} routes every ideal to exactly one of Sections~\ref{sec:p8}--\ref{sec:p11} or the analytic three-plane case. The three-plane case already gives
\(W_4\ge0\). On the remaining branches their conclusions give
\(D_0\ge m-29/10\).
This bound is positive, and \(A\ge m+1\). Hence \eqref{eq:p:1} gives
\[
\begin{aligned}
mW_4
&\ge(m+1)\left(m-\frac{29}{10}\right)-m(m-1)\\
&=-\frac{9m+29}{10}>-m,
\end{aligned}
\]
where the strict inequality is equivalent to \(m>29\). Thus \(W_4>-1\),
and its integrality gives \(W_4\ge0\). More generally, the same calculation
would accept a uniform deficit \(D_0\ge m-C\) for all \(m\ge30\) whenever
\(C<90/31\); the chosen value \(C=29/10\) lies just below this threshold.
Together with the three-plane case and the small-multiplicity argument, this
proves the theorem.
\end{proof}

\section{Analytic details}\label{sec:analytic}

The following three arguments are integral to the proof.

\subsection{Central boxes and three monotone horns}\label{sec:analytic-a}

\subsubsection{Statement}

For \(r\in\mathbb N\), write \([r]=\{0,1,\ldots,r\}\). For
\(z\in\mathbb N^3\), let

\[
\mathcal B(z)=[z_1]\times[z_2]\times[z_3].
\]

\begin{theorem}[Central-box/three-horn decomposition]\label{theorem:a:1.1}

Let \(U\subseteq\mathbb N^3\) be a nonempty finite lower ideal. Suppose
that no minimal point of \(\mathbb N^3\setminus U\) has all three
coordinates positive. Then there is a point \(c\in U\) and a disjoint
partition

\begin{equation}
U=\mathcal B(c)\mathbin{\dot\cup}A_1
 \mathbin{\dot\cup}A_2\mathbin{\dot\cup}A_3
\label{eq:a:1}
\end{equation}

with the following properties. For each \(i\), let \(j,k\) be the other
two coordinates, and put \(h_i=\max\{x_i:x\in U\}\), which is also the
largest coordinate-\(i\) axis level in \(U\). Then \(h_i\ge c_i\), and
there are nonincreasing functions

\begin{equation}
r_i:\{c_i+1,\ldots,h_i\}\longrightarrow\{0,\ldots,c_j\},
\qquad
s_i:\{c_i+1,\ldots,h_i\}\longrightarrow\{0,\ldots,c_k\}
\label{eq:a:2}
\end{equation}

such that

\begin{equation}
A_i=
\left\{x\in\mathbb N^3:
c_i<x_i\le h_i,\quad
0\le x_j\le r_i(x_i),\quad
0\le x_k\le s_i(x_i)
\right\}.
\label{eq:a:3}
\end{equation}

The domain in \eqref{eq:a:2} is empty when \(h_i=c_i\), giving an empty horn. Thus
the theorem includes a single box, one or two nonempty horns, zero central
caps, and ideals that do not contain all coordinate units.

\end{theorem}
\subsubsection{Plane determination and the compatibility graph}

The exclusion hypothesis has the following equivalent form:

\begin{equation}
x\in U
\quad\Longleftrightarrow\quad
(x_1,x_2,0),(x_1,0,x_3),(0,x_2,x_3)\in U.
\label{eq:a:4}
\end{equation}

The forward implication is lower closure. Conversely, if \(x\notin U\),
then \(x\) dominates a minimal excluded point \(p\). At least one coordinate
of \(p\) is zero, so one of the three planar points in \eqref{eq:a:4} also dominates
\(p\) and lies outside \(U\).

For coordinate \(i\), let

\[
h_i^{\mathrm{ax}}=\max\{t\in\mathbb N:te_i\in U\}.
\]

Construct a graph \(G\) with vertices

\[
V_i=\{(i,t):0\le t\le h_i^{\mathrm{ax}}\},
\qquad i=1,2,3.
\]

Each \(V_i\) is a clique. For \(i\ne j\), join \((i,t)\) to \((j,s)\)
exactly when \(te_i+se_j\in U\). The three zero-level vertices are
universal. Moreover, lower closure gives the nested-neighborhood property

\begin{equation}
t'<t\quad\Longrightarrow\quad
N(i,t)\setminus V_i\subseteq N(i,t')\setminus V_i.
\label{eq:a:5}
\end{equation}

\begin{lemma}\label{lemma:a:2.1}

The graph \(G\) is chordal.

\end{lemma}
\begin{proof}

Suppose an induced cycle has length at least four. Because there are only
three coordinate classes, it repeats a class. An induced cycle can contain at
most two vertices from a clique, and those two must be consecutive. Write
their levels as \(t'<t\). The other cycle neighbor of
\((i,t)\) lies outside \(V_i\); by \eqref{eq:a:5}, it is also adjacent to \((i,t')\).
This is a chord, a contradiction.
\end{proof}

\subsubsection{The required clique-tree lemma}

\begin{lemma}[Clique tree]\label{lemma:a:3.1}

The distinct maximal cliques of a finite nonempty connected chordal graph are
the nodes of a tree such that the nodes containing any fixed graph vertex
form a connected subtree.

\end{lemma}
\begin{proof}

We induct on the number of graph vertices. A complete graph has a one-node
tree. Otherwise choose nonadjacent vertices \(a,b\) and an inclusion-minimal
set \(Q\) separating them. Let \(C_a,C_b\) be their components in
\(G-Q\).

The separator \(Q\) is a clique. Indeed, minimality implies that every
\(q\in Q\) has a neighbor in both \(C_a\) and \(C_b\). If nonadjacent
\(q,r\in Q\) existed, take a shortest \(q\)-to-\(r\) path whose internal
vertices lie in \(C_a\), and another through \(C_b\). Their union would be
an induced cycle of length at least four: shortestness rules out chords
within either path, and distinct components of \(G-Q\) have no edges between
them. This contradicts chordality.

Let \(C_1,\ldots,C_d\), with \(d\ge2\), be all components of \(G-Q\).
Each induced graph on \(C_\ell\cup Q\) is connected, chordal, and smaller.
Apply induction to each, and choose in each resulting clique tree a bag
\(K_\ell\) containing the clique \(Q\). Join the nodes \(K_\ell\) in a
star. A vertex outside \(Q\) occurs in only one component tree. A vertex in
\(Q\) occurs in a connected subtree of every component tree containing
\(K_\ell\), and the added star joins those subtrees. Thus the resulting tree
of clique bags has the required connectedness property.

Every maximal clique of \(G\) occurs among these bags, because a clique cannot
meet two components of \(G-Q\). Some bags may be nonmaximal in \(G\), or the
same maximal clique may occur more than once. If a bag \(B\) is contained in
another bag \(C\), then every node on the tree path from \(B\) to \(C\)
contains every vertex of \(B\). Contract the first edge of that path and keep
the neighboring bag. This preserves the connectedness property. Repeating
removes all nonmaximal and duplicate bags, leaving exactly the distinct
maximal cliques.
\end{proof}

This finite clique-tree property is classical; see Gavril \cite{G}.

\subsubsection{From maximal cliques to maximal boxes}

Every maximal clique of the compatibility graph contains a prefix of each
coordinate class. Indeed, if it contains \((i,t)\), then every lower level in
\(V_i\) is adjacent to the whole clique by \eqref{eq:a:5}, and maximality adds it.
The universal zero levels ensure that no coordinate class is absent.

Let \(z_i\) be the largest level from \(V_i\) in the clique. All three pairs
of coordinate maxima are adjacent, so \eqref{eq:a:4} gives \(z=(z_1,z_2,z_3)\in U\).
Conversely, every \(z\in U\) gives the clique consisting of the three
coordinate prefixes through \(z_i\). Inclusion of these prefix cliques is
exactly coordinatewise comparison of their upper corners. It follows that
the maximal cliques of \(G\) correspond bijectively to the maximal points
of \(U\), and hence to the anchored boxes \(\mathcal B(z)\) maximal in
\(U\): a nonmaximal corner extends its prefix clique, and any extension of a
prefix clique produces a larger corner through \eqref{eq:a:4}.

Apply Lemma~\ref{lemma:a:3.1}, and write \(\mathcal T\) for the resulting tree on maximal
points. A graph vertex \((i,t)\) belongs to the clique for \(z\) precisely
when \(z_i\ge t\). Therefore

\begin{equation}
\boxed{
\{z\in\operatorname{Max}(U):z_i\ge t\}
\text{ induces a connected subtree whenever it is nonempty.}
}
\label{eq:a:6}
\end{equation}

\subsubsection{The median and the three arms}

Distinct maximal points of \(U\) are incomparable. Let \(z(v)\) denote the
point labeling a node \(v\) of \(\mathcal T\).

If \(v\) is a leaf with neighbor \(w\), choose a coordinate \(i\) for which
\(z_i(v)>z_i(w)\). If another node \(u\) had \(z_i(u)\ge z_i(v)\), the
superlevel subtree \eqref{eq:a:6} at level \(z_i(v)\) would contain \(u,v\) but not
\(w\), even though the path from \(u\) to \(v\) passes through \(w\).
Hence \(v\) is the unique global maximizer of coordinate \(i\). Two leaves
cannot be unique maximizers of the same coordinate, so \(\mathcal T\) has at
most three leaves.

For each coordinate \(i\), choose a node \(v_i\) maximizing \(z_i\). Every
leaf is one of \(v_1,v_2,v_3\), by the preceding uniqueness. Consequently
the whole tree is the minimal subtree spanning those three nodes: otherwise
a component outside that spanning subtree would end in an unselected leaf.

Let \(v_0\) be the median of \(v_1,v_2,v_3\), the unique common node of
their three pairwise paths, and set \(c=z(v_0)\). Repeated selected nodes are
allowed. The nontrivial paths from \(v_0\) to the \(v_i\) cover the tree and
meet only at \(v_0\); these are the at most three arms.

Consider an edge \(uv\) on the path from \(v_0\) toward \(v_i\), oriented
away from \(v_0\). If \(z_i(u)>z_i(v)\), the coordinate-\(i\) superlevel
subtree at height \(z_i(u)\) would contain \(u,v_i\) but not the intervening
node \(v\). Thus

\begin{equation}
z_i(u)\le z_i(v).
\label{eq:a:7}
\end{equation}

For \(j\ne i\), the selected maximizer \(v_j\) is either \(v_0\) or lies
on another arm. Hence \(u\) lies on the path from \(v\) to \(v_j\). If
\(z_j(v)>z_j(u)\), the corresponding superlevel subtree would contain
\(v,v_j\) but not \(u\). Therefore

\begin{equation}
z_j(u)\ge z_j(v),\qquad j\ne i.
\label{eq:a:8}
\end{equation}

Equality in \eqref{eq:a:7}, together with \eqref{eq:a:8}, would give \(z(v)\le z(u)\), contrary
to incomparability. Thus the outward coordinate increases strictly along
its arm, while both transverse coordinates weakly decrease.

\subsubsection{Disjoint rectangular sections}

All maximal boxes on arm \(i\) have transverse caps at most \(c_j,c_k\).
Fix an integer \(c_i<t\le h_i\). The central node and the other arms have
coordinate-\(i\) cap at most \(c_i\), so they contribute nothing at this
level. Among nodes on arm \(i\) whose outward cap is at least \(t\), let
\(v(t)\) be the first. All later contributing boxes have transverse
rectangles contained in that of \(v(t)\), by \eqref{eq:a:8}. Hence the entire level
section is

\begin{equation}
\{t\}\times\prod_{\ell\ne i}[z_\ell(v(t))],
\label{eq:a:9}
\end{equation}

after putting coordinate \(i\) first. As \(t\) increases, \(v(t)\) moves
weakly outward, so both transverse caps are nonincreasing. There is no gap
between \(c_i+1\) and the last nonempty level, because every anchored box
contains all lower levels.

Every point of \(U\) belongs to a maximal anchored box. If it is outside
\(\mathcal B(c)\), that maximal box lies on a unique arm, and the point
exceeds the corresponding central cap while staying below the other two.
It therefore lies in exactly one section \eqref{eq:a:9}. This proves the disjoint
partition \eqref{eq:a:1} and Theorem~\ref{theorem:a:1.1}.
\hfill\qedsymbol

Positive diagonal scaling preserves the decomposition. Rectangular
thickening replaces each integer level by a slab and each integer transverse
rectangle by a continuous rectangle; its two transverse side lengths remain
nonincreasing. Thus the thickened sets of Section~\ref{sec:p4} belong to the same
central-box/three-horn geometric class, up to null boundary conventions.

\subsubsection{One full-support corner}

Let \(T\subseteq\mathbb N^3\) be a finite lower ideal with exactly one
full-support minimal excluded point \(p\). Let \(P\) be the set of all
minimal excluded points of \(T\), and define

\begin{equation}
U=\mathbb N^3\setminus
\bigcup_{q\in P\setminus\{p\}}(q+\mathbb N^3).
\label{eq:a:10}
\end{equation}

\begin{corollary}[One-corner clipping]\label{corollary:a:7.1}

The set \(U\) is a finite lower ideal with no full-support minimal exclusion,
and

\begin{equation}
\boxed{T=U\setminus(p+\mathbb N^3).}
\label{eq:a:11}
\end{equation}

The ideals \(T\) and \(U\) have the same points on every coordinate axis.
If \eqref{eq:a:1} is a central-box/three-horn partition of \(U\), then

\begin{equation}
T=
\bigl(\mathcal B(c)\setminus(p+\mathbb N^3)\bigr)
\mathbin{\dot\cup}
\bigl(A_1\setminus(p+\mathbb N^3)\bigr)
\mathbin{\dot\cup}
\bigl(A_2\setminus(p+\mathbb N^3)\bigr)
\mathbin{\dot\cup}
\bigl(A_3\setminus(p+\mathbb N^3)\bigr).
\label{eq:a:12}
\end{equation}

At an outward level \(t\) of horn \(i\), with transverse coordinates
\(j,k\) and uncut caps \(r,s\), the retained section is the exact set

\begin{equation}
\begin{cases}
[r]\times[s], & t<p_i,\\
([r]\times[s])\setminus
   \bigl([p_j,r]\times[p_k,s]\bigr), & t\ge p_i,
\end{cases}
\label{eq:a:13}
\end{equation}

where an interval \([a,r]\) is empty when \(a>r\).

\end{corollary}
\begin{proof}

Every member of \(P\setminus\{p\}\) has support at most two, so the minimal
excluded points of \(U\), after redundant generators are removed, also have
support at most two. For each coordinate \(i\), the first missing axis point
of \(T\) is a pure-axis member of \(P\). It remains in \eqref{eq:a:10}, proving that
\(U\) is finite. The full-support orthant \(p+\mathbb N^3\) contains no axis
point, so \(T\) and \(U\) have the same axis sections. The upper-ideal
description of a lower ideal's complement gives \eqref{eq:a:11}.

Deleting the same upper orthant from each piece of a disjoint partition
preserves disjointness and proves \eqref{eq:a:12}. At level \(t<p_i\), no point lies
above \(p\); at level \(t\ge p_i\), precisely the transverse upper-right
rectangle lies above \(p\). This proves \eqref{eq:a:13}.
\end{proof}

The same clipping formulas remain valid for any positive integer point
\(p\), even if it is inactive or not minimal excluded after clipping. Allowing
such choices enlarges a later optimization class and therefore remains safe
for an upper-bound certificate. In particular, if no point of \(U\) dominates
\(p\), the clipping is inactive and the representation includes the
no-full-corner case.

\subsubsection{The nested-horn recurrence}

\begin{lemma}[Nested-horn recurrence]\label{lemma:horn-recurrence}

Fix an outward coordinate \(i\), a first level \(t_0=c_i+1\), and central
transverse caps \(b_0=c_j\), \(d_0=c_k\). Suppose a problem assigns an
additive score \(C_i(t,r,s)\) to an uncut rectangular section or the retained
section \eqref{eq:a:13}, and has a pointwise feasibility predicate
\(\mathcal F_i(t,r,s)\). Let
\(F_i(t,b,d)\) be the maximum score of a possibly empty nested sequence
starting at level \(t\), with first caps at most \(b,d\). Then

\begin{equation}
\boxed{
F_i(t,b,d)=\max\left\{
0,
\max_{\substack{0\le r\le b,\ 0\le s\le d\\
                 \mathcal F_i(t,r,s)}}
\bigl(C_i(t,r,s)+F_i(t+1,r,s)\bigr)
\right\}.
}
\label{eq:a:14}
\end{equation}

A terminal level beyond the allowed axis range has value zero.
\end{lemma}
\begin{proof}

The terminal value is the induction base case. The zero option terminates
the horn. Every nonempty uncut section contains its axis
point, which survives clipping because \(p\) is positive; lower closure then
forbids a later nonempty section after termination. Every other choice in
\eqref{eq:a:14} selects the current rectangle and passes exactly its caps to the next
level. Induction on the remaining number of levels proves that \eqref{eq:a:14} covers
every feasible nested horn and only such sequences. There are finitely many
levels and cap choices, so backward evaluation terminates.
\end{proof}

Equivalently, a prefix-maximum implementation may compare the state
\((b,d)\) with \((b-1,d)\) and \((b,d-1)\), because every proper
subrectangle decreases at least one cap.

For a fixed center, the three arms in \eqref{eq:a:1} or \eqref{eq:a:12} are disjoint, so their
additive optima may be computed independently and added to the central-piece
score. Maximizing over all centers therefore covers every ideal in the
theorem or clipping corollary. Particular score formulas, degree or weighted
height allowances, and interval bounds are specified in Section~\ref{sec:proof}.


\subsection{Full weighted Apéry ideals}\label{sec:analytic-b}


\subsubsection{Setup and statement}

Let

\[
S=\langle m,a_1,a_2,a_3\rangle
\]

be a minimally four-generated numerical semigroup. Let
\(T\subseteq\mathbb N^3\) be its preferred Apéry lower ideal, and put

\[
\lambda(x)=a_1x_1+a_2x_2+a_3x_3,
\qquad
M=\max_{x\in T}\lambda(x).
\]

Thus \(|T|=m\), the labels of its points are pairwise distinct modulo
\(m\), and \(0,e_1,e_2,e_3\in T\). We say that \(T\) is \emph{full weighted} if

\begin{equation}
T=\{x\in\mathbb N^3:\lambda(x)\le M\}.
\label{eq:b:1}
\end{equation}

Permuting the three coordinates preserves \eqref{eq:b:1}, lower closure, residue
injectivity, and all projection quantities used below. We may therefore
write the three nonmultiplicity generators as

\[
a<b<d
\]

and use coordinates \((x,y,z)\), so that

\[
\lambda(x,y,z)=ax+by+dz.
\]

\begin{theorem}[Full-weighted-ideal theorem]\label{theorem:b:1.1}

If the preferred Apéry ideal \(T\) is full weighted, then

\[
W_4(S)\ge0.
\]

\end{theorem}
\subsubsection{Classification of the column projection}

Project away the \(x\)-coordinate. By \eqref{eq:b:1}, its image is

\begin{equation}
Q=\pi_x(T)
=\{(y,z)\in\mathbb N^2:by+dz\le M\}.
\label{eq:b:2}
\end{equation}

Put

\[
h_y=\left\lfloor\frac Mb\right\rfloor,
\qquad
h_z=\left\lfloor\frac Md\right\rfloor.
\]

Both are positive because \(e_2,e_3\in T\).

\begin{lemma}[Three possible projections]\label{lemma:b:2.1}

The set \(Q\) is one of

\begin{equation}
\begin{aligned}
Q_3&=\{(0,0),(1,0),(0,1)\},\\
Q_4&=\{(0,0),(1,0),(2,0),(0,1)\},\\
Q_6&=\{(y,z)\in\mathbb N^2:y+z\le2\}.
\end{aligned}
\label{eq:b:3}
\end{equation}

\end{lemma}
\begin{proof}

For \(1\le j\le h_y\), define

\begin{equation}
p_j=
\left(
\left\lfloor\frac{M-bj}{a}\right\rfloor+1,
j,0
\right).
\label{eq:b:4}
\end{equation}

This is a minimal excluded point. Its \(x\)-predecessor is included by the
floor definition. Moreover,

\[
\lambda(p_j)\le M+a,
\]

so its \(y\)-predecessor has weight at most \(M+a-b\le M\). Both displayed
coordinates of \(p_j\) are positive because \(bj\le M\).

The representative in \(T\) of the residue of \(p_j\) must have support
disjoint from \(\{x,y\}\), so it lies on the \(z\)-axis. The corners
\(p_1,\ldots,p_{h_y}\) have intersecting supports and hence distinct
residues. The \(z\)-axis contains exactly

\[
0,e_3,\ldots,h_ze_3,
\]

and therefore

\begin{equation}
h_y\le h_z+1.
\label{eq:b:5}
\end{equation}

First suppose \(M<b+d\). Then \eqref{eq:b:2} contains no point with both coordinates
positive. Since

\[
d\le M<b+d<2d,
\]

we have \(h_z=1\). Equation \eqref{eq:b:5} gives \(h_y\le2\), while \(h_y\ge1\).
This gives \(Q_3\) or \(Q_4\).

Now suppose \(M\ge b+d\). The point

\begin{equation}
p=
\left(
\left\lfloor\frac{M-b-d}{a}\right\rfloor+1,
1,1
\right)
\label{eq:b:6}
\end{equation}

is a full-support minimal exclusion: its \(x\)-predecessor is included by
the floor definition, while removing \(e_2\) or \(e_3\) reduces its weight
from at most \(M+a\) by at least \(a\). By the corner rules, \(p\) has
residue zero. Each \(p_j\) has support intersecting that of \(p\), so none
of their distinct representatives on the \(z\)-axis can be zero. Hence

\begin{equation}
h_y\le h_z.
\label{eq:b:7}
\end{equation}

If \(M\ge2b+d\), then

\[
p'=
\left(
\left\lfloor\frac{M-2b-d}{a}\right\rfloor+1,
2,1
\right)
\]

is a second full-support minimal exclusion. Indeed, its \(x\)-predecessor is
included by the floor definition, and removing \(e_2\) or \(e_3\) reduces
its weight from at most \(M+a\) by at least \(a\). It is distinct from
\(p\) because its \(y\)-coordinate is two. This contradicts the uniqueness
of a full-support corner, so

\[
M<2b+d<3d
\]

and \(h_z\le2\). On the other hand,

\[
M\ge b+d>2b,
\]

so \(h_y\ge2\). Since \(b<d\) also gives \(h_y\ge h_z\), equation \eqref{eq:b:7}
forces

\begin{equation}
h_y=h_z=2.
\label{eq:b:8}
\end{equation}

It follows from \eqref{eq:b:8} that \(2d\le M<3b\). Every \((y,z)\) of total degree at
most two has weight at most \(2d\), while every point of total degree at
least three has weight at least \(3b\). Thus \eqref{eq:b:2} is exactly \(Q_6\).

\end{proof}

\subsubsection{The six-column triangle is impossible}

For \((y,z)\in Q\), let

\begin{equation}
L_{yz}=\left\lfloor\frac{M-by-dz}{a}\right\rfloor+1
\label{eq:b:9}
\end{equation}

be the length of its \(x\)-column. Subtracting \(b\) or \(d\) from the
available weight decreases a positive column length by at least one, because
\(b,d>a\).

\begin{lemma}[Six-column exclusion]\label{lemma:b:3.1}

The projection \(Q_6\) in \eqref{eq:b:3} cannot occur.

\end{lemma}
\begin{proof}

Suppose it does, and abbreviate its positive column lengths by

\begin{equation}
\begin{array}{c|cccccc}
(y,z)&(0,0)&(1,0)&(2,0)&(0,1)&(1,1)&(0,2)\\ \hline
L_{yz}&\alpha&\beta&\gamma&\varphi&\varepsilon&\eta.
\end{array}
\label{eq:b:10}
\end{equation}

The strict column descents give

\begin{equation}
\alpha>\beta>\gamma>0,
\qquad
\alpha>\varphi>\eta>0,
\qquad
\varepsilon<\min(\beta,\varphi).
\label{eq:b:11}
\end{equation}

Work in \(\mathbb Z/m\mathbb Z\), and let \(r\) be the least positive
integer for which \(ra\equiv0\pmod m\). The point
\((\varepsilon,1,1)\) is the full-support minimal exclusion above the mixed
column: its \(x\)-predecessor lies in that column, and its other two
predecessors are included by \(\varepsilon<\varphi,\beta\). It therefore has
residue zero:

\begin{equation}
\varepsilon a+b+d\equiv0\pmod m.
\label{eq:b:12}
\end{equation}

The residues in the origin column are

\[
0,a,\ldots,(\alpha-1)a,
\]

whereas \eqref{eq:b:12} makes those in the mixed column

\[
-\varepsilon a,-(\varepsilon-1)a,\ldots,-a.
\]

All \(\alpha+\varepsilon\) points belong to \(T\), so residue injectivity
shows that these consecutive multiples of \(a\) are distinct. Consequently

\begin{equation}
r\ge\alpha+\varepsilon.
\label{eq:b:13}
\end{equation}

The two points \((\beta,1,0)\) and \((\gamma,2,0)\) are minimal
exclusions: their \(x\)-predecessors lie at the tops of their columns, and
their \(y\)-predecessors are included by \(\beta<\alpha\) and
\(\gamma<\beta\). Their representatives have support on the \(z\)-axis.
Their residues are nonzero because their supports meet that of
\((\varepsilon,1,1)\), and they are distinct from each other. Since the
\(z\)-axis consists of \(0,e_3,2e_3\), their residues therefore permute
\(d,2d\).

The swapped assignment would be

\begin{equation}
\beta a+b\equiv2d,
\qquad
\gamma a+2b\equiv d.
\label{eq:b:14}
\end{equation}

Subtracting twice the second congruence in \eqref{eq:b:14} from the first gives

\[
(\beta-2\gamma)a-3b\equiv0\pmod m,
\]

while the second congruence in \eqref{eq:b:14}, substituted into \eqref{eq:b:12}, gives

\[
(\varepsilon+\gamma)a+3b\equiv0\pmod m.
\]

Adding these relations eliminates \(b\) and gives

\[
(\beta-\gamma+\varepsilon)a\equiv0\pmod m.
\]

But

\[
0<\beta-\gamma+\varepsilon
<\alpha+\varepsilon\le r,
\]

contradicting the definition of \(r\). We must therefore have

\begin{equation}
\beta a+b\equiv d,
\qquad
\gamma a+2b\equiv2d.
\label{eq:b:15}
\end{equation}

The same argument for the minimal exclusions
\((\varphi,0,1)\) and \((\eta,0,2)\), whose representatives lie on the
\(y\)-axis, rules out the swapped assignment there: it would give
\((\varphi-\eta+\varepsilon)a\equiv0\pmod m\), although

\[
0<\varphi-\eta+\varepsilon<\alpha+\varepsilon\le r.
\]

Consequently

\begin{equation}
\varphi a+d\equiv b,
\qquad
\eta a+2d\equiv2b.
\label{eq:b:16}
\end{equation}

Adding the first congruences in \eqref{eq:b:15}--\eqref{eq:b:16}, and subtracting the second
congruences from twice their respective first congruences, shows that

\begin{equation}
\beta+\varphi,
\qquad
2\beta-\gamma,
\qquad
2\varphi-\eta
\label{eq:b:17}
\end{equation}

are positive multiples of \(r\). Each is strictly less than \(2\alpha\),
and \eqref{eq:b:13} gives \(\alpha<r\). Thus all three integers in \eqref{eq:b:17} lie strictly
between zero and \(2r\), so

\[
\beta+\varphi=r,
\qquad
2\beta-\gamma=r,
\qquad
2\varphi-\eta=r.
\]

Adding the last two equalities and subtracting twice the first gives

\[
\gamma+\eta=0,
\]

contrary to \eqref{eq:b:11}. Hence \(Q_6\) is impossible.
\end{proof}

\begin{remark}\label{remark:b:3.2}

The proof of Lemma~\ref{lemma:b:3.1} used only residue bijectivity, the corner rules, the
six-point projection, and the strict descents in \eqref{eq:b:11}. The exact floor
formula \eqref{eq:b:9} is not otherwise needed. This isolates the arithmetic obstruction
from the full-weighted classification which produces the six-point shape.

\end{remark}
\subsubsection{The remaining projection scores}

By Section~\ref{sec:p3}, it suffices to prove \(\Phi(T,K)\ge0\), where
\(K=\operatorname{Max}(T)\). The two remaining projections admit one
calculation.

\begin{lemma}[Three- and four-column projections]\label{lemma:b:4.1}
If \(Q=Q_k\), \(k\in\{3,4\}\), then \(\Phi(T,K)\ge0\).
\end{lemma}
\begin{proof}
Let \(\alpha\) be the origin-column length, and index the other lengths
in the order displayed in \eqref{eq:b:3}. Since \(b,d>a\), the lengths
decrease strictly along each axis of the projection:
\begin{center}
\begin{tabularx}{\textwidth}{@{}lYY@{}}
\toprule
Projection & Column lengths & Strict descents \\
\midrule
\(Q_3\) & \(\alpha,\beta,\gamma\) & \(\alpha>\beta,\gamma\ge1\) \\
\(Q_4\) & \(\alpha,\beta,\gamma,\delta\) &
\(\alpha>\beta>\gamma\ge1,\ \alpha>\delta\ge1\) \\
\bottomrule
\end{tabularx}
\end{center}
These descents imply \(\alpha\ge k-1\) and make the column tops pairwise
incomparable, so they are precisely \(K\).

Every non-origin column lies on exactly one axis of the projection. The
two projections retaining the column direction therefore count each
column once and the origin column twice. Hence
\begin{equation}
|K|=k,\qquad P(T)=k+\alpha+m.
\label{eq:b:20}
\end{equation}
For these two shapes the sum of the projected total degrees is
\(2k-4\), and the sum of their coordinate maxima is \(k-1\). Thus
\[
\sum_{x\in K}|x|_1=m+k-4,\qquad
\sum_j\max_{x\in K}x_j=\alpha+k-2.
\]
It follows that
\begin{equation}
E(K)=m-\alpha-2,\qquad
\Phi(T,K)=2(\alpha-k+1)\ge0.
\label{eq:b:21}
\end{equation}
\end{proof}

\subsubsection{Proof of the full-weighted theorem}

\begin{proof}[Proof of Theorem~\ref{theorem:b:1.1}]

Lemma~\ref{lemma:b:2.1} lists the three possible \(x\)-column projections. Lemma~\ref{lemma:b:3.1}
excludes the six-column triangle, and Lemma~\ref{lemma:b:4.1} gives
\(\Phi(T,K)\ge0\) in the remaining cases. The projection inequality
\eqref{eq:p:6} proves \(W_4(S)\ge0\).

\end{proof}


\subsection{The continuous no-full-corner moment inequality}\label{sec:analytic-c}


\subsubsection{Box-horn sets and their moment functional}

Boundary faces never affect the integrals below. A \emph{box-horn set of height
one} is the disjoint union, up to null sets, of a central box

\begin{equation}
C=[0,a]\times[0,b]\times[0,c],
\qquad a,b,c\ge0,\qquad a+b+c\le1,
\label{eq:c:5}
\end{equation}

and at most three outward horns. The horn beyond the \(a\)-face has the form

\begin{equation}
\{(t,y,z):a<t<1,\ 0<y<f(t),\ 0<z<g(t)\},
\label{eq:c:6}
\end{equation}

where \(f,g\) are nonnegative, nonincreasing finite step functions satisfying

\begin{equation}
f(t)\le b,\qquad g(t)\le c,
\qquad t+f(t)+g(t)\le1.
\label{eq:c:7}
\end{equation}

The other two horns are defined by coordinate permutation. Empty horns and
zero-width sections are allowed.

For a measurable \(E\subseteq\mathbb R_{\ge0}^3\), define

\begin{equation}
\mathcal J(E)=\int_E\bigl(3(x_1+x_2+x_3)-2\bigr)\,dx.
\label{eq:c:8}
\end{equation}

Thus \(\mathcal J(E)\le0\) is equivalent to

\begin{equation}
\frac1{\operatorname{vol}E}
\int_E(x_1+x_2+x_3)\,dx\le\frac23
\label{eq:c:9}
\end{equation}

when \(E\) has positive volume. The central box contributes

\begin{equation}
\mathcal J(C)
=abc\left(\frac32(a+b+c)-2\right).
\label{eq:c:10}
\end{equation}

The contribution of the horn \eqref{eq:c:6} is

\begin{equation}
\mathcal J_a(f,g)
=\int_a^1 f(t)g(t)
 \left(3t+\frac32(f(t)+g(t))-2\right)dt.
\label{eq:c:11}
\end{equation}

We next obtain a sharp enough common envelope for \eqref{eq:c:11}, then show that the
three envelopes are absorbed by the negative contribution \eqref{eq:c:10}.

\subsubsection{The boundary-tail potential}

It is convenient to give the formulas at an arbitrary height allowance
\(L>0\). At this height write
\[
\mathcal J_L(E)=\int_E\bigl(3(x_1+x_2+x_3)-2L\bigr)\,dx,
\]
so that \(\mathcal J=\mathcal J_1\). For \(0\le x\le y\), define

\begin{equation}
\begin{aligned}
P_L(x,y)={}&\frac{Lx^3}{6}-\frac{x^4}{4}
 +\frac{Lxy^2}{2}-\frac{3x^2y^2}{4}-\frac{xy^3}{2},\\
g(q,z)={}&\frac{q^2z^2}{8}+\frac{q^3z}{24}+\frac{q^4}{192},\\
Q_L(x,y)={}&P_L(x,y)
 +g\bigl((y+2x-L)_+,y-x\bigr).
\end{aligned}
\label{eq:c:12}
\end{equation}

Extend \(Q_L\) symmetrically in \(x,y\).

\begin{lemma}[Boundary-tail optimum]\label{lemma:c:3.1}

Suppose a horn is already on the height boundary: in reverse time
\(r=L-t\), its transverse sides are nondecreasing functions with sum \(r\),
ending at caps \(x\le y\). Its contribution is at most \(Q_L(x,y)\).

\end{lemma}
\begin{proof}

By scaling, take \(L=1\). At each \(r\), sort the two sides and write their
difference as \(d(r)\ge0\). Sorting preserves the sum, product, caps, and
monotonicity. The boundary contribution becomes

\begin{equation}
\int_0^{x+y}\frac{(2-3r)(r^2-d(r)^2)}8\,dr.
\label{eq:c:13}
\end{equation}

The two sides are nondecreasing precisely when \(d\) is \(1\)-Lipschitz.
The cap conditions give

\begin{equation}
\max(0,r-2x)\le d(r)\le\min(r,2y-r).
\label{eq:c:14}
\end{equation}

The coefficient of \(-d^2\) changes sign only at \(r=2/3\). After fixing
\(d(2/3)\), the optimal admissible function is therefore the smallest one
allowed by \eqref{eq:c:14} to the left of \(2/3\), and the largest one allowed to its
right. Writing \(q=(2/3-d(2/3))/2\), the associated transverse sides are

\begin{equation}
\begin{cases}
(r/2,r/2),&0\le r\le2q,\\
(q,r-q),&2q\le r\le y+q,\\
(r-y,y),&y+q\le r\le x+y.
\end{cases}
\label{eq:c:15}
\end{equation}

The same profile follows directly from pointwise product maximization if
\(x+y\le2/3\). When \(x+y\ge2/3\), feasibility gives

\[
\max(0,2/3-y)\le q\le\min(x,1/3).
\]

Let \(I(q;x,y)\) denote the integral in \eqref{eq:c:13} for \eqref{eq:c:15}. Direct integration
gives

\begin{equation}
\frac{\partial I}{\partial q}
=\frac{(y-q)^2(1-y-2q)}2.
\label{eq:c:16}
\end{equation}

Hence the optimum is attained at

\begin{equation}
q_*=\min\left(x,\frac{1-y}{2}\right).
\label{eq:c:17}
\end{equation}

In the second case above, \(y\ge1/3\), so \(q_*\) lies in the displayed
feasible interval. The endpoint \(q=x\) has value \(P_1(x,y)\). If
\(q_*<x\), integration of \eqref{eq:c:16} from \(q_*\) to \(x\) shows that the
improvement is

\begin{equation}
g(y+2x-1,y-x).
\label{eq:c:18}
\end{equation}

Equations \eqref{eq:c:12} and \eqref{eq:c:18} prove the claim for \(L=1\); homogeneity gives the
general statement.
\end{proof}

\subsubsection{Reducing an arbitrary horn to one slab}

For \(a+x+y\le L\), put

\begin{equation}
F_L(a;x,y)=Q_L(x,y)
 +\frac{3a-L}{2}xy(L-a-x-y),
\label{eq:c:19}
\end{equation}

and, whenever \(a+B+C\le L\), define

\begin{equation}
U_L(a;B,C)=
\max_{0\le x\le B,\ 0\le y\le C}F_L(a;x,y).
\label{eq:c:20}
\end{equation}

The term added to \(Q_L\) is the contribution of a constant initial slab
from \(t=a\) to \(t=L-x-y\).

\begin{lemma}[One-slab envelope]\label{lemma:c:4.1}

Every finite-step horn based at \(a\), with transverse caps \(B,C\) and
height allowance \(L\), has contribution at most \(U_L(a;B,C)\).

\end{lemma}
\begin{proof}

Scale to \(L=1\), and sort the two transverse coordinates so \(B\le C\).
Replacing the two sides pointwise by their minimum and maximum preserves
nestedness and area. The smaller side remains at most \(B\), while both
original sides and hence the larger remain at most \(C\). We may therefore
assume that the first side is the smaller one at every level.
We proceed in three stages: make the slabs height-tight, compare them
with a boundary interpolation, and remove or merge the remaining jumps.

\smallskip\noindent\emph{Step 1: endpoint tightness.}
First fix the finite list of nested rectangle states and vary only the
outward endpoints of their slabs. The objective is a separable convex
quadratic in those endpoints: the quadratic coefficient at an endpoint is
\(3(A_i-A_{i+1})/2\ge0\), where \(A_i\) is the area of the \(i\)-th rectangle
and \(A_{n+1}=0\). A maximum on the endpoint polytope is attained at a
vertex. A block of equal endpoints at such a vertex either begins at \(a\)
or meets one of its height bounds. Removing zero-length slabs leaves every
surviving slab tight at its right endpoint.

\smallskip\noindent\emph{Step 2: boundary interpolation.}
Reverse time by \(r=1-t\), and write

\[
A(r)=1-\frac32r.
\]

If a tight slab starts at reverse time \(r\), ends with rectangle \((x,y)\),
and \(x+y=s\le r\), its exact contribution is

\begin{equation}
A(r)(r-s)xy.
\label{eq:c:21}
\end{equation}

Retain the possibly slack initial slab, including any central slack;
the filling below applies only to transitions between actual tight states.
For a transition from reverse time \(r\) to \(s\), where
\(0\le s\le r\le2/3\), let \(q(\rho)\) be the area of a monotone boundary
interpolation and \(xy\) the smaller endpoint area. Then
\[
\int_s^r A(\rho)q(\rho)\,d\rho
\ge xy\int_s^r A(\rho)\,d\rho
\ge xyA(r)(r-s).
\]
The right side is the tight-slab contribution \eqref{eq:c:21}, so filling
such transitions weakly increases the objective.

If none
of the states reaches reverse time at most \(2/3\), all terms in \eqref{eq:c:21} are
nonpositive and the empty horn dominates. Otherwise let \((u,v)\), \(u\le v\),
be the first tight state with \(u+v\le2/3\).
From \((u,v)\) toward reverse time zero, use the pointwise maximizing
continuation: keep the smaller side \(u\) fixed down to \((u,u)\), then
follow the balanced path. This continuation is nested and respects both
caps; it is the boundary profile from Lemma~\ref{lemma:c:3.1} in this
nonnegative region.

It remains to control the earlier transitions between actual states, where
\(A(r)\le0\). At every such state of sum
\(s\), nestedness and the cap conditions allow its product to be replaced by
the smaller product at

\begin{equation}
x_*(s)=\max(u,s-C),\qquad
y_*(s)=\min(C,s-u).
\label{eq:c:22}
\end{equation}

Indeed \(x(s-x)\) is increasing for \(x\le s/2\); the pair in \eqref{eq:c:22} is
nested as \(s\) varies, respects \(B,C\), and ends at \((u,v)\). Since the
coefficient in \eqref{eq:c:21} is nonpositive, this replacement can only increase the
objective. All earlier states now lie on the L-shaped boundary path

\begin{equation}
(B,C)\longrightarrow(u,C)\longrightarrow(u,u).
\label{eq:c:23}
\end{equation}

The possibly slack initial position is represented only for bookkeeping by
the virtual state \((x_0,C)\), where \(x_0=1-a-C\ge B\). The first actual
jump covers the virtual excess.

\begin{figure}[htbp]
\centering
\setlength{\unitlength}{1pt}
\begin{picture}(270,160)
\thinlines
\put(15,15){\vector(1,0){230}}
\put(15,15){\vector(0,1){135}}
\put(255,13){\makebox(0,0){\(x\)}}
\put(13,157){\makebox(0,0){\(y\)}}
\linethickness{0.8pt}
\put(125,130){\line(-1,0){65}}
\put(93,130){\vector(-1,0){15}}
\put(60,130){\line(0,-1){70}}
\put(60,106){\vector(0,-1){15}}
\put(125,130){\circle*{4}}
\put(60,130){\circle*{4}}
\put(60,60){\circle*{4}}
\put(60,142){\makebox(0,0){\((u,C)\)}}
\put(125,142){\makebox(0,0){\((B,C)\)}}
\put(60,44){\makebox(0,0){\((u,u)\)}}
\linethickness{0.4pt}
\multiput(128,130)(6,0){13}{\line(1,0){3}}
\put(205,130){\circle{4}}
\put(205,142){\makebox(0,0){\((x_0,C)\)}}
\qbezier[18](95,130)(77.5,112.5)(60,95)
\put(95,130){\circle*{4}}
\put(60,95){\circle*{4}}
\put(110,99){\makebox(0,0)[l]{\footnotesize crossing jump}}
\end{picture}
\caption{Schematic of the L-shaped boundary path in \eqref{eq:c:23}. The dashed
extension represents the virtual initial state; a crossing jump skips
the turn.}
\label{fig:boundary-path}
\end{figure}

\smallskip\noindent\emph{Step 3: removing and merging jumps.}
We now compare a jump on \eqref{eq:c:23} with filling the same part of the boundary
path continuously. Along a branch with one fixed side \(c\), a jump of
length \(d\) has gain

\begin{equation}
G_c(d)=\frac c4d^2(2d+3c-2),
\label{eq:c:24}
\end{equation}

and two jumps satisfy

\begin{equation}
G_c(d+e)-G_c(d)-G_c(e)
=\frac{3c}{2}de\left(d+e+c-\frac23\right).
\label{eq:c:25}
\end{equation}

These follow by integrating \(c(r-c)A(r)\).

\smallskip\noindent\emph{Jumps after the turn.}
On the second branch of \eqref{eq:c:23},
the fixed side is \(u\), and every jump has
\(d\le C-u\le1-2u\). Thus \(2d+3u-2\le-u\), so \eqref{eq:c:24} is nonpositive and all
such jumps may be filled.

\smallskip\noindent\emph{A jump crossing the turn.}
At most one jump crosses the turn in \eqref{eq:c:23}, as illustrated in
Figure~\ref{fig:boundary-path}. Suppose it starts at
\((u+p,C)\) and ends at \((u,C-q)\). If its gain over the continuous L-path
is denoted by \(G(p,q;u,C)\), direct integration gives

\begin{equation}
\frac{\partial G}{\partial q}
=\frac u2(p+q)(3p+3q+3u-2).
\label{eq:c:26}
\end{equation}

This derivative changes sign at most once, from negative to positive, so the
maximum occurs at \(q=0\) or \(q=C-u\). The continuous baseline from the
jump's start through the remaining tail is independent of \(q\), so the same
endpoint conclusion applies to the whole continuation. The first choice
moves the jump entirely onto the fixed-\(C\) branch.

For the balanced endpoint, a jump from reverse time \(r\) to \((z,z)\),
followed by the balanced tail, has value

\[
E_r(z)=A(r)(r-2z)z^2+\frac23z^3-\frac32z^4,
\]

with

\begin{equation}
E_r'(z)=-6z\left(z-\frac r2\right)
              \left(z-r+\frac23\right).
\label{eq:c:27}
\end{equation}

On \(0\le z\le R\le r/2\), the maximum of \(E_r\) is therefore at \(0\) or
\(R\). For a noninitial crossing jump, take \(R\) to be the smaller side of
its actual starting state; the second endpoint is a fixed-smaller-side jump
and is dominated by filling. For the initial crossing jump, take \(R=B\),
which is valid because \(2B\le B+C\le1-a\). Its second endpoint is the
admissible initial slab ending at \((B,B)\). If the zero endpoint wins, every
earlier term still lies in the nonpositive part, so the empty horn dominates.
Thus the crossing jump either disappears or is replaced by an admissible
one-slab construction.

\smallskip\noindent\emph{Merging jumps before the turn.}
All remaining jumps lie on the fixed-\(C\) branch. Retain the first jump,
including any mandatory initial slack, and fill every later jump whose gain
in \eqref{eq:c:24} is nonpositive. If a later jump of length \(e\) has positive gain,
then either \(C\ge2/3\), or

\[
e>\frac32\left(\frac23-C\right)>\frac23-C.
\]

Equation \eqref{eq:c:25} shows that merging it into the first jump cannot decrease the
total gain. The individual gains depend only on jump lengths, so the jump and
continuous lengths may be reordered along the same branch before merging.
Their total length, the continuous baseline, nestedness, and the initial cap
constraint are unchanged. Repetition leaves a single initial jump followed
by a boundary profile.

If that jump ends at effective caps \((x,y)\), its contribution together with
the optimal continuation is \(F_1(a;x,y)\). The effective caps lie inside
\((B,C)\), so \eqref{eq:c:20} proves the lemma.
\end{proof}

\subsubsection{The joint three-horn inequality}

\begin{lemma}[Joint envelope]\label{lemma:c:5.1}

If \(a,b,c\ge0\), \(S=a+b+c\le L\), then

\begin{equation}
\boxed{
U_L(a;b,c)+U_L(b;a,c)+U_L(c;a,b)
\le abc\left(2L-\frac32S\right).
}
\label{eq:c:28}
\end{equation}

\end{lemma}
\begin{proof}

For \(x\le y\), first omit the boundary-tail gain from \eqref{eq:c:12} and write

\begin{equation}
H_L(a;x,y)=P_L(x,y)
 +\frac{3a-L}{2}xy(L-a-x-y).
\label{eq:c:29}
\end{equation}

Let \(r=L-a-x-y\). Direct differentiation gives

\begin{equation}
(H_L)_y=\frac{xr}{2}\,[3(a+y)-L].
\label{eq:c:30}
\end{equation}

The gain in \eqref{eq:c:12} is active only when \(y+2x>L\). Since \(x\le y\),
activity implies \(y>L/3\), and hence \(3(a+y)-L>0\). Thus
\eqref{eq:c:30} is nonnegative wherever the nondecreasing gain is active.
Before activation, the displayed derivative changes sign at most once, from
negative to positive. Therefore, for fixed \(x\), \(F_L(a;x,y)\) first
decreases and then increases as \(y\) runs from \(x\) to its allowed cap;
its maximum is at the diagonal or at the full larger cap. If
\(a+x\ge L/3\), \eqref{eq:c:30} sends the diagonal endpoint to the full
cap. If \(a+x\le L/3\), the gain vanishes on the diagonal segment and

\begin{equation}
\frac d{dt}H_L(a;t,t)
=-6t\left(t-\frac{L-a}{2}\right)
       \left(t-\left(\frac L3-a\right)\right)\le0
\label{eq:c:31}
\end{equation}

for \(0\le t\le x\). That diagonal endpoint is at most zero, which is
attained by the empty horn. Thus, for \(B\ge C\),

\begin{equation}
U_L(a;B,C)=\max_{0\le x\le C}F_L(a;x,B).
\label{eq:c:32}
\end{equation}

Relabel the central sides so that \(a\le b\le c\), and first take the
saturated case \(L=S\). The gain in \eqref{eq:c:12} vanishes for the horns based at
\(b\) and \(c\). For the gain-free expression, one also has

\begin{equation}
2(H_L)_x=(a-x)(x-y)^2
 +r\,[x^2+2ay-yr],
\label{eq:c:33}
\end{equation}

where \(a\) in this formula denotes the horn's base and \(y\) its full larger
cap. Equation \eqref{eq:c:33} is nonnegative in the two relevant ranges. Consequently
the \(b\)- and \(c\)-horns take their full smaller cap \(a\). Only the horn
based at the smallest side remains to be optimized.

Indeed, for the \(b\)-horn one has \(0\le x\le a\le b\), larger cap \(c\),
and \(r=a-x\); for the \(c\)-horn the larger cap is \(b\) and the same
relations hold. Both terms on the right of \eqref{eq:c:33} are then nonnegative.

For \(0\le x\le b\), set

\begin{equation}
\Delta=\frac{abcS}{2}-F_S(a;x,c)
       -H_S(b;a,c)-H_S(c;a,b).
\label{eq:c:34}
\end{equation}

Three nonnegative polynomial identities cover its entire domain. In the
gain-free range \(x\le a\), substitute

\[
x=p,\quad a=p+q,\quad b=p+q+r,\quad c=p+q+r+s.
\]

Then

\begin{equation}
\begin{aligned}
\Delta={}&\frac34p^2(q+r)^2
 +\frac12p(q+r)^2(s+r+2q)\\
&+q^2\left[\frac{(s+2r)^2}{4}
 +\frac23q(s+2r)+\frac12q^2\right]\ge0.
\end{aligned}
\label{eq:c:35}
\end{equation}

In the remaining gain-free range \(a\le x\le(a+b)/2\), substitute

\[
a=p,\quad x=p+q,\quad b=p+q+r,\quad c=p+q+r+s.
\]

Then

\begin{equation}
\begin{aligned}
\Delta={}&\frac34p^2r^2+\frac12pr^2(s+r+2q)
 +\frac12qr^2(s+r)\\
&+q^2\left[\frac{(s+2r)^2}{4}+\frac{r^2}{4}\right]
 +\frac13q^3(s+2r)+\frac16q^4\ge0.
\end{aligned}
\label{eq:c:36}
\end{equation}

Finally, the gain is active exactly when \(x\ge(a+b)/2\). Substitute

\[
a=p,\quad x=p+2q+r,\quad b=p+2q+2r,
\quad c=p+2q+2r+s.
\]

After including the gain in \(F_S\), identity \eqref{eq:c:34} becomes

\begin{equation}
\begin{aligned}
\Delta={}&\frac34p^2r^2+\frac12pr^2(s+3r+4q)\\
&+s^2\left(\frac14r^2+qr+\frac12q^2\right)\\
&+s\left(\frac{11}{6}r^3+7qr^2+7q^2r+\frac73q^3\right)\\
&+\frac{31}{12}r^4+\frac{34}{3}qr^3
 +\frac{33}{2}q^2r^2+\frac{31}{3}q^3r
 +\frac{31}{12}q^4\ge0.
\end{aligned}
\label{eq:c:37}
\end{equation}

All variables in \eqref{eq:c:35}--\eqref{eq:c:37} are nonnegative. These cases prove \eqref{eq:c:28} when
\(L=S\).

It remains to allow central slack. Fix effective smaller caps

\[
0\le x\le b,\qquad0\le y,z\le a,
\]

for the horns based at \(a,b,c\), respectively; by \eqref{eq:c:32}, their larger caps
are \(c,c,b\). First omit the possible gain and define

\begin{equation}
\begin{aligned}
G(L)={}&abc\left(2L-\frac32S\right)
 -H_L(a;x,c)-H_L(b;y,c)-H_L(c;z,b).
\end{aligned}
\label{eq:c:38}
\end{equation}

For \(L=S+\delta\), direct expansion gives

\begin{equation}
G(S+\delta)=G(S)+\delta\Lambda
 +\frac{\delta^2}{2}(cx+cy+bz),
\label{eq:c:39}
\end{equation}

where

\begin{equation}
\begin{aligned}
\Lambda={}&2abc-f_a(x)-f_b(y)-f_c(z),\\
f_a(t)={}&\frac{t^3}{6}+\frac{ct^2}{2}+c(a-b)t,\\
f_b(t)={}&\frac{t^3}{6}+\frac{ct^2}{2}+c(b-a)t,\\
f_c(t)={}&\frac{t^3}{6}+\frac{bt^2}{2}+b(c-a)t.
\end{aligned}
\label{eq:c:40}
\end{equation}

These three functions are convex on their intervals; \(f_b,f_c\) are
nondecreasing, while \(f_a\) is maximized at \(0\) or \(b\). At those two
endpoint choices, respectively,

\[
\Lambda\ge
a^2\left(\frac{b+c}{2}-\frac a3\right)\ge0,
\]

\begin{equation}
\Lambda\ge
\frac{(b-a)^2[2(b-a)+3(c-b)]}{6}\ge0.
\label{eq:c:41}
\end{equation}

Thus the gain-free deficit is nondecreasing in \(L\). At \(L=S\), it is
nonnegative for every \(x,y,z\): equations \eqref{eq:c:35}--\eqref{eq:c:37} handle the largest
possible \(b\)- and \(c\)-horns, whose full smaller cap is \(a\).

Only the horn based at \(a\) can have the extra gain from \eqref{eq:c:12}, because
\(c+2y\le S\) and \(b+2z\le S\) for the other two. Its gain

\[
g((c+2x-L)_+,c-x)
\]

is nonincreasing in \(L\). Subtracting it preserves the monotonicity of the
full deficit. Since the saturated deficit is nonnegative by \eqref{eq:c:35}--\eqref{eq:c:37}, it
remains nonnegative for \(L\ge S\). Maximizing the three horns independently
proves \eqref{eq:c:28}.
\end{proof}

\begin{theorem}[Continuous no-full-corner inequality]\label{theorem:c:5.2}

Every positive-volume box-horn set \(K\) of height one satisfies

\begin{equation}
\boxed{
\frac1{\operatorname{vol}K}
\int_K(x_1+x_2+x_3)\,dx\le\frac23.
}
\label{eq:c:42}
\end{equation}

\end{theorem}
\begin{proof}

Apply Lemma~\ref{lemma:c:4.1} to the three horns and then Lemma~\ref{lemma:c:5.1} with \(L=1\).
Their total contribution to \eqref{eq:c:8} is at most

\[
abc\left(2-\frac32(a+b+c)\right),
\]

which is the negative of the central contribution \eqref{eq:c:10}. Hence
\(\mathcal J(K)\le0\), which is \eqref{eq:c:42}.
\end{proof}

\begin{remark}
The constant \(2/3\) is sharp for finite-step box-horn sets as a supremum.
The central cube \([0,1/3]^3\), with three square-section horns of side
\((1-t)/2\) at height \(1/3<t<1\), has volume \(1/9\) and coordinate-sum
moment \(2/27\). Approximating the horns from below by nested finite-step
sections preserves the height bound and makes the normalized first moment
tend to \(2/3\). Thus a uniformly stronger bound requires structure beyond
the box-horn hypotheses.
\end{remark}


\subsubsection{Application to the thickening}

The rectangular thickening of Section~\ref{sec:p4}, after positive diagonal
scaling, has the box-horn form proved in Section~\ref{sec:analytic-a}. Its simplex containment
and the preceding theorem give its deficit at least one third.
Here \(a,b,c\) denote box side lengths.

\subsubsection{The degree-four cardinality bound}


\begin{lemma}\label{lemma:c:7.1}

If \(T\) has no full-support minimal exclusion and every point of \(T\) has
total degree at most four, then

\begin{equation}
|T|\le29.
\label{eq:c:44}
\end{equation}

\end{lemma}
\begin{proof}

Use the Section~\ref{sec:analytic-a} decomposition with central point
\(q=(q_1,q_2,q_3)\). Necessarily \(|q|_1\le4\). At horn \(i\), level
\(t>q_i\), its rectangular section has caps \(r\le q_j\), \(s\le q_k\) and
\(t+r+s\le4\). Consequently

\begin{equation}
\begin{aligned}
|T|\le{}&\prod_{i=1}^3(q_i+1)\\
&+\sum_{i=1}^3\sum_{t=q_i+1}^4
\max_{\substack{0\le r\le q_j,\ 0\le s\le q_k\\t+r+s\le4}}
(r+1)(s+1).
\end{aligned}
\label{eq:c:45}
\end{equation}

Dropping the nesting relation between successive sections only enlarges this
upper bound. The expression is symmetric in \(q\), so the possible sorted
central triples and their values in \eqref{eq:c:45} are:

\begin{center}
\small
\begin{tabular}{@{}llll@{}}
\toprule
\(q\) & Bound in \eqref{eq:c:45} & \(q\) & Bound in \eqref{eq:c:45} \\
\midrule
\((0,0,0)\) & 13 & \((0,0,1)\) & 19 \\
\((0,0,2)\) & 23 & \((0,0,3)\) & 25 \\
\((0,0,4)\) & 25 & \((0,1,1)\) & 25 \\
\((0,1,2)\) & 28 & \((0,1,3)\) & 28 \\
\((0,2,2)\) & 28 & \((1,1,1)\) & 29 \\
\((1,1,2)\) & 29 &  &  \\
\bottomrule
\end{tabular}
\end{center}

These are all partitions into at most three parts of an integer at most four.
The largest entry is \(29\), proving \eqref{eq:c:44}.
\end{proof}

\section{Exact finite verification}\label{sec:verification}

The proof has seven retained local finite premises. Their mathematical
contracts appear in Section~\ref{sec:proof}; the accompanying packet
contains two checking paths for each:

\begin{center}
\small
\begin{tabularx}{\textwidth}{@{}lYY@{}}
\toprule
Premise & First path & Second path \\
\midrule
Generator box & Cyclic residue distances & Ordinary membership \\
No-corner interval & Prefix horn maxima & Direct recursive transitions \\
Short corner, high & Subtractive moments/prefix maxima & Disjoint moments/direct transitions \\
Short corner, low & Recursive profiles/ten offsets & Cartesian profiles/successors \\
Degree-six profiles & Recursive profiles/bitsets & Subset profiles/columns \\
Three-allowance strip & Subtractive moments/prefix maxima & Retained points/all transitions \\
Residual high interval & Full-box subtraction & Disjoint-box statistics \\
\bottomrule
\end{tabularx}
\end{center}

The three-plane case has no finite premise. The published theorem for
\(m\le19\) is an external input, not a component of this replay.
The two profile paths generate their complete families independently,
including all stated degree, cardinality, drop, and surface tests.
No saved dual, exceptional-family list, LP solver, modular lift, or
floating-point tolerance is an acceptance premise. The two inline
short-corner witnesses are reconstructed and checked directly.

\subsection{Closed interval certificates}

The three interval inputs store only a root and indexed tree topology
at scale \(q=4096\). Their closed roots have height ranges
\(5\ldots24\), \(6\ldots42\), and \(7\ldots78\), respectively,
with each weight coordinate ranging from one to the root's maximum
height. At every node, the checker applies the stated tightening,
then either derives two children sharing a strictly interior split
wall or checks the declared leaf predicate. Node zero is the root;
every node must be reached exactly once. Hence no component, split
boundary, or subregion may be omitted.

Empty leaves require an inverted tightened enclosure. Analytic leaves
are permitted only in the first two trees, under the rules given in
Sections~\ref{sec:p8} and~\ref{sec:p9}. Every other nonempty leaf
requires a fresh complete evaluation of the horn recurrence. Each
short-corner leaf checks all three orientations. Every residual leaf
checks the full Cartesian corner range from Section~\ref{sec:p11}; an
empty range is explicitly reported as vacuity. The two coverage walks
use depth-first integer and breadth-first rational arithmetic,
respectively. Neither walk trusts cached bounds or stored child boxes.

For reference, the verified trees have the following node counts:
\begin{center}
\small
\begin{tabular}{@{}lrrrr@{}}
\toprule
Tree & Splits & DP leaves & Analytic leaves & Empty leaves \\
\midrule
No corner & 25,442 & 24,912 & 531 & 0 \\
Short corner & 55,432 & 54,978 & 455 & 0 \\
Residual high & 47,229 & 47,088 & 0 & 142 \\
\bottomrule
\end{tabular}
\end{center}
These counts are diagnostic, not substitutes for coverage or leaf tests.
The residual replay checks 35,852,138 corner bounds; its largest DP
score is 11,824, below \(29\cdot4096/10\).
The strip contains 715 sorted corner/allowance configurations across
allowances 18, 19, and 20. The omitted null translation requires no
allowance-21 computation.

\subsection{Arithmetic and replay}

% BEGIN selected PRELIM arithmetic
All finite predicates are exact integer comparisons. Python integers are
unbounded. C++ evaluators require signed scores with at least 63 value bits,
indices with at least 31 value bits, and 64-bit bitsets where used.
For the largest interval root, \(N_i\le78\), \(q=4096\), a conservative
bound on score/moment/threshold intermediates is
\[
16\cdot79^3\bigl(156(3\cdot78q)+235q\bigr)<2^{51}<2^{63}-1.
\]
The strip envelope \(8\cdot21^3(9\cdot20+4)<2^{31}\) is smaller.
Degree-six profiles have \(m\le84\), \(s_i\le504\), axis caps at most six;
their denominator-cleared margins fit below \(2^{30}\) in magnitude.
Generator-box distances are at most \((m-1)m(m-2)\le21924\).
The concrete implementations also check their required integer widths.
No mathematical operation is performed on a negative-infinity sentinel.
% END selected PRELIM arithmetic

The packet is the repository directory \texttt{prelim/}. From that
directory, run \texttt{python3 -I -B verify.py}. Only Python standard
libraries and a C++17 compiler are required; Git, network access,
historical files, and third-party Python packages are not used.
The runner hashes every manifest input before and after execution,
compiles fresh evaluators with
\texttt{-std=c++17 -O3 -Wall -Wextra -pedantic}, and rejects compiler
diagnostics, partial runs, unsupported modes, and malformed responses.
The two checking paths keep separate mathematical state and results.
Scheduling and administrative hashing are not shared evaluators.

A complete isolated replay of the frozen packet passed all seven pairs
and 35 rejection tests. The generated record \texttt{replay.json}
identifies tested inputs, toolchains, coverage, per-leaf agreement,
and predicate partitions. The manifest and replay record identify
specific sources; the replay record is never an acceptance input.
The interval data total approximately 3.4 MB and contain no cached
leaf answers. The TeX sources compile independently of that packet;
the packet is needed to reproduce the computational premises.

\section{Extensions to embedding dimension greater than four}
\label{sec:extensions}

The Ap\'ery framework extends to all embedding dimensions. The geometric
surplus argument does not presently do so. We separate these two facts:
this section gives identities and a conditional transfer principle, but
does not claim Wilf's conjecture for any unrestricted \(e>4\).

\subsection{Dimension-independent moment and collision identities}

Let \(e=d+1\), \(S=\langle m,a_1,\ldots,a_d\rangle\), and define
\(A,w,H,T,s\) as before, with \(T\subseteq\mathbb N^d\). Preferred
factorizations again give \(|T|=m\), lower closure, and distinct residues.
For the indexed family of coordinate lines, the same progression sums give
\[
\sum_\ell L_\ell=dm,\qquad
\sum_\ell L_\ell t_\ell=(d+1)s.
\]
Consequently Zhai's weighted downset baseline \cite{Z} reads
\[
D_e:=dmH-(d+1)w\cdot s\ge0.
\]
The genus formula and \(W_e=en-c=(e-1)c-eg\) give the exact identity
\begin{equation}\label{eq:extension-moment}
mW_e=A D_e-\frac{e-2}{2}\,m(m-1).
\end{equation}
The factor \(m(m-1)\) makes the correction integral even when its
coefficient is a half-integer.

The predecessor-subtraction collision proof also remains valid. A minimal
excluded point has support disjoint from that of its representative; two
distinct minimal exclusions with intersecting support have distinct
residues. In particular there is still at most one full-support corner,
with representative zero. Unlike in three exponent coordinates, the
remaining corners can have support three or more.

Finally, if nonnegative masses on \(T\) have total \(dm\) and moment
\((d+1)s+\beta\), with \(\beta\ge0\), their destinations have weight
at most \(H\), and \(w_i\ge1\), then
\[
D_e=\sum_x\alpha_x(H-w\cdot x)+w\cdot\beta\ge|\beta|_1.
\]
Thus the local centroid mechanism extends. What is missing is a uniform
surplus of the required size for every relevant ideal.

\subsection{A conditional mesh-transfer principle}

The translated-mesh argument has a useful dimension-independent form.
Let \(K\) be a measurable positive-volume downset in the unit simplex.
Choose an integer \(N>d\), put \(h=1/N\), and define
\(T_z=\{y\in\mathbb N^d:z+hy\in K\}\).
Suppose that, for almost every \(z\in[0,h)^d\), this induced discrete
ideal satisfies
\[
d r_z|T_z|-(d+1)\sum_{y\in T_z}|y|_1\ge\alpha|T_z|,
\qquad
r_z=\left\lfloor\frac{1-|z|_1}{h}\right\rfloor.
\]
Only allowances \(N-d,\ldots,N-1\) occur almost everywhere. Writing
\(\epsilon_z=1-|z|_1-hr_z\), the integrand expands as
\[
\begin{aligned}
\sum_{y\in T_z}\bigl(d-(d+1)|z+hy|_1\bigr)
&=hD_{r_z}(T_z)+(d\epsilon_z-|z|_1)|T_z|\\
&\ge(\alpha h-|z|_1)|T_z|.
\end{aligned}
\]
After integrating, each coordinate remainder has mean at most \(h/2\)
by the interval-fiber calculation in Section~\ref{sec:p11}. Hence
\begin{equation}\label{eq:extension-mesh}
d-(d+1)\mathbb E_K|X|_1
\ge\left(\alpha-\frac d2\right)\frac1N.
\end{equation}
Our strip transfer is \(d=3\), \(N=21\), \(\alpha=4\), giving \(5/42\).
In higher dimension this is a conditional tool: the required finite
inequality and an exhaustive way to check it have not been established
here. A positive gap needs \(\alpha>d/2\).

\subsection{Why the three-horn decomposition does not extend directly}

The compatibility graph is chordal in three coordinate classes because
an induced cycle of length at least four repeats a class. With four
classes, that argument fails, and the failure is real.

Consider the finite lower ideal in \(\mathbb N^4\)
\[
U=\{x\in\{0,1\}^4:x_1x_3=0,\ x_2x_4=0\}.
\]
Its minimal exclusions are \(2e_i\) for \(1\le i\le4\), and
\(e_1+e_3,e_2+e_4\). None has full support. The compatibility graph
on the four positive axis levels has the induced cycle
\[
(1,1)\;\text{---}\;(2,1)\;\text{---}\;(3,1)\;\text{---}\;(4,1)\;\text{---}\;(1,1),
\]
whose two diagonals are absent.

More directly, no central box \([0,c]\), \(c\in U\), together with
four horns restricted to exceeding only one central coordinate can
cover \(U\). The support of \(c\) is empty, a singleton, or an adjacent
pair of this cycle. In every case there is an adjacent pair disjoint
from that support. Its vector \(e_i+e_j\) lies in \(U\), exceeds two
central caps, and belongs to neither the central box nor any such horn.

This is an obstruction for arbitrary lower ideals, not an assertion
that \(U\) admits an Ap\'ery labeling. A higher-dimensional proof might
exploit additional arithmetic restrictions or replace the decomposition
with a richer structural model. The current three-horn recurrence cannot
simply be reused with more axes.

Known dimension-independent results still cover substantial regions,
including \(3e\ge m\) \cite{Graph}, \(c\le3m\) \cite{Macaulay}, and
\(t\le e-1\), where \(t\) is the type \cite{FGH}. They do not supply
the missing uniform geometric surplus for a fixed \(e>4\).

\section{Related work}\label{sec:related}

\subsection{Weighted ideals, L-shapes, and graph structure}

Zhai's asymptotic theorem \cite{Z} establishes an approximate Wilf
inequality for fixed embedding dimension outside finitely many semigroups.
Our argument instead seeks a uniform, exact result in embedding dimension
four. Within the weighted-ideal framework of Zhai and
Hellus--Rechenauer--Waldi \cite{HRW}, the additional task is a quantitative
surplus sufficient to absorb the genus correction, rather than a new
construction of the exponent ideal.

Aguil\'o-Gost, Garc\'ia-S\'anchez, and Llena \cite{L} study the number
of L-shapes for semigroups of embedding dimension four. Chomicz
\cite{C} gives a geometric procedure for constructing Ap\'ery sets
in that dimension and applies it to explicit families. This is closely
related geometric work, but its relation-deletion and rearrangement
results are not premises of our finite families. We retain the canonical
lexicographic representatives. For any choice of one factorization per
Ap\'ery element, the sum of labels is fixed; rearranging representatives
alone cannot improve that sum. We therefore introduce no optimization
or realizability filter justified only by a change of L-shape.

The clique-tree theorem used in our central-box decomposition is classical
\cite{G}. The application to three-coordinate lower ideals and the resulting
moment bounds are proved explicitly in Section~\ref{sec:analytic-a}.
This use of graph theory is different from Eliahou's graph-theoretic
criterion \(3e\ge m\) \cite{Graph}.

\subsection{Known regions and fixed-multiplicity computation}

Fr\"oberg, Gottlieb, and H\"aggkvist \cite{FGH} prove \(g\le tn\).
Thus \(t\le e-1\) implies
Wilf immediately. Sammartano \cite{Sammartano} proves the criterion
\(2e\ge m\); Eliahou \cite{Graph} improves it to \(3e\ge m\).
Eliahou's Macaulay-theorem argument \cite{Macaulay} settles
\(c\le3m\). These results describe important previously settled regions,
but are not additional assumptions in our branch partition.

The computations of Bruns et al. \cite{B} and Kliem--Stump \cite{KS}
use Kunz cones and their faces. Our small-multiplicity computation instead
enumerates the bounded generator box obtained from the counterexample
reduction in Section~\ref{sec:p6}. This is a different finite assertion,
not a reproduction of the published Kunz-cone computations.

\subsection{Type and presentation bounds}

Chomicz's subsequent paper \cite{CT} bounds the cardinality \(\eta(S)\)
of minimal presentations in terms of the type:
\[
t(S)-11\le\eta(S)\le4t(S)+5.
\]
These bounds concern another aspect of the same Ap\'ery geometry.
They do not give a uniform bound \(t\le3\) for four generators and
therefore do not replace our large-multiplicity branches.
Marashdeh \cite{M} derives type bounds involving Ap\'ery extremal
statistics and reduces Wilf's conjecture to a further inequality.
We do not use those reductions as established substitutes for the
quantitative moment surplus required here.

The cited 2026 works are preprints in the versions listed in the
bibliography. The literature assessment is dated September 17, 2026;
we make no claim that this bibliography certifies publication priority.

\section{How the proof was constructed}\label{sec:construction}

Karthik Sethuraman began this project through several weeks of iteration
with GPT-6 Astra. This produced an initial complete candidate proof and
associated exploratory computations.

Sethuraman then iterated with GPT-5.6 Sol via Codex to decompose, simplify,
and rewrite the proof. This work organized the argument into explicit proof
obligations and dependencies, reconstructed the Ap\'ery foundations,
separated the geometric branches, normalized notation and endpoints, and
reduced the scope of finite verification. Several computations were replaced
by analytic proofs. In particular, the three-plane branch became entirely
analytic, the low short-corner branch was reduced to local centroid predicates
and two explicit witnesses, and the degree-six branch used complete
local-or-axis predicates. The mesh transfer was reduced to the three
allowances that occur almost everywhere.

The reconstruction also included a literature review that identified and
attributed the preferred-factorization, weighted-downset, initial-ideal, and
support ingredients, and situated the proof relative to Chomicz's Ap\'ery-set
constructions. The resulting comparison clarified what survives in more
exponent coordinates and where the three-dimensional structure fails.

Computer assistance has two distinct roles here. AI assisted discovery,
proof development, exposition, and implementation. Exact programs establish
the finite assertions specified mathematically in the paper, as described in
Section~\ref{sec:verification}.

\bibliographystyle{alpha}
\bibliography{references}
\end{document}
